% ============================================================================
%  中国农业周期性分析报告 —— 第一部分（农业子行业周期性分析）单独成册
%  编译：xelatex main-part1.tex （两遍 + biber），或 make part1
%  内容 = 封面 + 摘要（第一部分口径）+ 目录 + 图表目录 + 第一章至第六章 + 参考文献
%  正文内容与 main.tex 完全同源（\input 同一批 tex/ch0*.tex），不复制、不改写；
%  仅前置部分与收录范围不同。
% ============================================================================
\documentclass[11pt, a4paper, fontset=none, zihao=-4]{ctexart}
% ============================================================================
%  preamble.tex —— 中国农业周期性分析报告 导言区
%  编译：xelatex（需 ctex + tikz/pgfplots + Source Han 字体）
% ============================================================================

% ---------- 中文与字体 ----------
% 注意：ctexart 类已加载 ctex，此处不可重复 \usepackage{ctex}
% （重复加载会触发 ctexhook Warning 并破坏 fontspec 参数解析）
\usepackage{fontspec}
\setmainfont[AutoFakeBold=2.5]{Source Han Serif CN}
\setsansfont[AutoFakeBold=2.5]{Source Han Sans CN}
\setmonofont{WenQuanYi Zen Hei Mono}
\setCJKmainfont[AutoFakeBold=2.5]{Source Han Serif CN}
\setCJKsansfont[AutoFakeBold=2.5]{Source Han Sans CN}
\setCJKmonofont{WenQuanYi Zen Hei Mono}
\newCJKfontfamily\heiti{Source Han Sans CN}
\newCJKfontfamily\songti{Source Han Serif CN}
\newCJKfontfamily\heiB[AutoFakeBold=5]{Source Han Sans CN}

% ---------- 版面 ----------
\usepackage[top=2.6cm, bottom=2.6cm, left=2.6cm, right=2.6cm]{geometry}
\usepackage{setspace}
\onehalfspacing
\sloppy
\setlength{\parskip}{2pt plus 1pt}
\setlength{\parindent}{2em}
\usepackage{fancyhdr}
\usepackage{titlesec}

% ---------- 数学与符号 ----------
\usepackage{amsmath, amssymb, amsfonts, mathtools}
\usepackage{bm}

% ---------- 表格 ----------
\usepackage{booktabs}
\usepackage{longtable}
\usepackage{array}
\usepackage{tabularx}
\usepackage{multirow}
\usepackage{makecell}
\usepackage{threeparttable}
\usepackage{siunitx}
\sisetup{group-separator={,}, group-minimum-digits=4, detect-weight=true}

% ---------- 图形 ----------
\usepackage{graphicx}
\usepackage{xcolor}
\usepackage{tikz}
\usepackage{pgfplots}
\usepackage{pgfplotstable}
\pgfplotsset{compat=1.18}
\usepgfplotslibrary{fillbetween, groupplots, statistics, patchplots}
\usetikzlibrary{arrows.meta, positioning, shapes.geometric, calc, fit,
                backgrounds, patterns, decorations.pathreplacing, shadows.blur,
                matrix, chains, fadings, intersections}
% TikZ 外部库（如未安装 external 依赖则自动忽略，不影响编译）
\usetikzlibrary{external}

% ---------- 参考文献（gb7714-2015 国标，biber 后端） ----------
\usepackage[backend=biber, style=gb7714-2015, sorting=none,
            url=true, doi=false, isbn=false, eprint=false]{biblatex}
\addbibresource{refs.bib}
\setlength{\bibitemsep}{2pt}
\renewcommand*{\bibfont}{\footnotesize}

% ---------- 颜色主题（农本绿） ----------
\definecolor{AgriGreen}{HTML}{2E6E4E}
\definecolor{AgriLight}{HTML}{6FA287}
\definecolor{AgriPale}{HTML}{E4EFE8}
\definecolor{AgriGold}{HTML}{C7A252}
\definecolor{HogPink}{HTML}{C97B84}
\definecolor{Grain}{HTML}{D9A441}
\definecolor{OilBlue}{HTML}{3B6E9E}
\definecolor{SoftRed}{HTML}{B5533C}
\definecolor{AquaTeal}{HTML}{2C8C87}
\definecolor{InkGray}{HTML}{3A3F44}
\definecolor{LightGray}{HTML}{F2F3F4}

% 供 pgfplots 循环使用的品种配色
\pgfplotscreateplotcyclelist{agri}{
  {AgriGreen, thick},
  {HogPink, thick},
  {Grain, thick},
  {OilBlue, thick},
  {SoftRed, thick},
  {AquaTeal, thick},
  {AgriGold, thick},
  {InkGray, thick},
  {AgriLight, thick},
  {violet, thick},
}

% 通用折线图样式
\pgfplotsset{
  % 年份坐标轴：x 为十进制年份，刻度标签取整
  yearticks/.style={
    xticklabel style={/pgf/number format/fixed,
                      /pgf/number format/precision=0,
                      /pgf/number format/1000 sep={},
                      font=\footnotesize},
    x tick label as interval=false,
  },
  % 相关系数条形图（y 轴为类别编号）
  catbar/.style={
    ytick=data, y dir=reverse,
    yticklabel style={font=\scriptsize},
    xlabel style={font=\small},
    tick label style={font=\footnotesize},
    axis line style={InkGray, line width=0.5pt},
    grid=major, grid style={LightGray, line width=0.3pt},
    enlarge y limits=0.03,
  },
  agri axis/.style={
    width=13.2cm, height=6.2cm,
    tick label style={font=\footnotesize},
    label style={font=\small},
    title style={font=\small\heiti},
    legend style={font=\scriptsize, fill=white, fill opacity=0.85,
                  text opacity=1, draw=AgriLight!50, rounded corners=1pt},
    grid=major, grid style={LightGray, line width=0.3pt},
    axis line style={InkGray, line width=0.5pt},
    every axis plot/.append style={line width=1pt},
  },
}

% ---------- 列表、盒子 ----------
\usepackage{enumitem}
\setlist[itemize]{itemsep=2pt, topsep=3pt, parsep=0pt, leftmargin=1.6em}
\setlist[enumerate]{itemsep=2pt, topsep=3pt, parsep=0pt, leftmargin=1.8em}
\usepackage[most]{tcolorbox}

\newtcolorbox{keybox}[1][核心结论]{
  enhanced, breakable, colback=AgriPale, colframe=AgriGreen,
  boxrule=0.4pt, left=6pt, right=6pt, top=4pt, bottom=4pt,
  fonttitle=\sffamily\bfseries\small, title={#1},
  attach boxed title to top left={xshift=6pt, yshift=-2.4pt},
  boxed title style={colback=AgriGreen, colframe=AgriGreen, size=small},
  arc=1.5pt, before skip=6pt, after skip=6pt}

\newtcolorbox{noteh}[1][数据说明]{
  enhanced, breakable, colback=LightGray, colframe=InkGray!40,
  boxrule=0.3pt, left=6pt, right=6pt, top=4pt, bottom=4pt,
  fonttitle=\sffamily\bfseries\scriptsize, title={#1},
  arc=1pt, before skip=5pt, after skip=5pt}

% ---------- 标题格式 ----------
\ctexset{
  section/format = {\heiti\zihao{4}\color{AgriGreen}},
  subsection/format = {\heiti\zihao{-4}\color{InkGray}},
  subsubsection/format = {\heiti\zihao{-4}\color{InkGray!80!black}},
  section/aftername = \quad,
  subsection/aftername = \quad,
}
% 部（Part）样式：ctexart 无 chapter，故用 titlesec 定制 \part
\titleclass{\part}{straight}[\section]
\titleformat{\part}[display]
  {\centering\heiti\zihao{2}\color{AgriGreen}}
  {\zihao{5}\sffamily\color{AgriGold} 第\,\zhnumber{\thepart}\,部分}
  {12pt}
  {}
  [\vspace{-8pt}\textcolor{AgriLight}{\rule{0.6\textwidth}{0.8pt}}]
\titlespacing*{\part}{0pt}{0pt}{28pt}
\usepackage{zhnumber}
\titleformat{\paragraph}[runin]{\heiti\zihao{-4}}{}{0pt}{}[：]

% ---------- 代码与超链接 ----------
\usepackage{listings}
\lstset{basicstyle=\ttfamily\scriptsize, breaklines=true,
        frame=single, rulecolor=\color{InkGray!40},
        backgroundcolor=\color{LightGray}, columns=fullflexible}
\usepackage[colorlinks=true, linkcolor=AgriGreen, citecolor=AgriGreen,
            urlcolor=OilBlue, bookmarksnumbered=true,
            CJKbookmarks=true]{hyperref}
\usepackage{footmisc}
\usepackage{caption}
\captionsetup{font={small,sf}, labelfont={bf,color=AgriGreen},
              labelsep=quad, skip=4pt}
\captionsetup[table]{position=top}

% ---------- 自定义宏 ----------
% 图表数据来源脚注（放在图表正下方）
\newcommand{\srcfile}[1]{\par\vspace{-2pt}%
  {\footnotesize\sffamily\color{InkGray!85}资料来源：#1\par}}
% 数据口径说明
\newcommand{\datacal}[1]{{\footnotesize\sffamily\color{InkGray!75}（#1）}}
% 周期阶段标记
\newcommand{\phase}[1]{\textcolor{AgriGreen}{\textbf{【#1】}}}
% 关键数字高亮
\newcommand{\hl}[1]{\textcolor{SoftRed}{\textbf{#1}}}
\newcommand{\up}{\textcolor{SoftRed}{$\blacktriangle$}}
\newcommand{\down}{\textcolor{AgriGreen}{$\blacktriangledown$}}
% 公司代码
\newcommand{\stk}[1]{\texttt{#1}}
% 缺失数据标记
\newcommand{\nadata}{\textcolor{InkGray!60}{---}}


\title{中国农业周期性分析报告（第一部分）}
\author{农业行业研究}
\date{2026 年 9 月}

\begin{document}

% ---------------------------------------------------------------- 封面
% ---------------------------------------------------------------- 封面
\begin{titlepage}
\begin{tikzpicture}[remember picture, overlay]
  % 顶部色带
  \fill[AgriGreen] (current page.north west)
    rectangle ([yshift=-4.2cm]current page.north east);
  \fill[AgriLight!55] ([yshift=-4.2cm]current page.north west)
    rectangle ([yshift=-4.45cm]current page.north east);
  \fill[AgriGold] ([yshift=-4.45cm]current page.north west)
    rectangle ([yshift=-4.62cm]current page.north east);
  % 底部色带
  \fill[AgriPale] (current page.south west)
    rectangle ([yshift=2.6cm]current page.south east);
  \fill[AgriGreen] ([yshift=2.6cm]current page.south west)
    rectangle ([yshift=2.72cm]current page.south east);
  % 稻穗装饰（右下水印）
  \begin{scope}[shift={([xshift=-3.4cm, yshift=8.2cm]current page.south east)},
                scale=1.05, opacity=0.18]
    \foreach \a in {-58,-48,-38,-28,-18,-8}{
      \fill[AgriGreen] (0,0) .. controls (0.35,1.0) and (0.55,1.8) .. (0.18,2.5)
        .. controls (-0.05,1.9) and (-0.2,1.0) .. (0,0);
      \draw[AgriGreen, line width=0.7pt] (0,0) -- ({1.4*cos(\a)},{1.4*sin(\a)});
    }
    \draw[AgriGreen, line width=1.6pt] (0,-0.4) -- (0,2.7);
    \foreach \a in {-52,-40,-28}{
      \draw[AgriGreen, line width=0.9pt] ({-0.15},{1.0+0.55*\a/50}) -- ({1.1*cos(\a)},{1.0+0.55*\a/50+1.1*sin(\a)});
    }
  \end{scope}
  % 标题文字
  \node[anchor=west, text=white, align=left]
    at ([xshift=2.2cm, yshift=-2.05cm]current page.north west)
    {\heiti\fontsize{26}{32}\selectfont 中国农业周期性分析报告};
  \node[anchor=west, text=AgriPale, align=left]
    at ([xshift=2.2cm, yshift=-3.35cm]current page.north west)
    {\heiti\fontsize{12.5}{16}\selectfont 基于申万行业分类的品种周期机制与上市公司画像};
  \node[anchor=west, text=AgriPale!85, align=left]
    at ([xshift=2.25cm, yshift=-3.95cm]current page.north west)
    {\sffamily\fontsize{9.5}{12}\selectfont
     China Agriculture: Cycle Analysis and Listed-Company Profiles};
\end{tikzpicture}

\vspace*{4.2cm}

\begin{center}
\begin{tikzpicture}
  \node[draw=AgriGreen, line width=0.8pt, fill=AgriPale, rounded corners=2pt,
        inner sep=12pt, text width=12.6cm, align=left]
  {\small\sffamily
   \textbf{研究范围}\quad 申万行业分类（2021 版）农林牧渔一级行业及其下属
   二级、三级子行业（种植业、养殖业、饲料、动物保健、农产品加工、渔业、林业）。\par\vspace{4pt}
   \textbf{第一部分}\quad 主要农产品品种的供需周期机制分析；以猪周期为基准，
   考察其他养殖品种（肉牛、肉羊、禽、水产）与经济作物的周期相关性与传导路径。\par\vspace{4pt}
   \textbf{第二部分}\quad 近三年农业全部 A 股上市公司的公司画像：基本情况、
   周期敏感度、盈利能力、财务结构、战略转型、历史融投资与未来潜在融资需求。\par\vspace{4pt}
   \textbf{数据来源}\quad akshare 开源数据库、巨潮资讯网、新浪财经、东方财富、
   农业农村部及国家统计局公开数据，截至 2026 年 9 月。};
\end{tikzpicture}
\end{center}

\vfill
\begin{center}
  {\heiti\zihao{4}\color{AgriGreen} 农业行业研究}\\[4pt]
  {\sffamily\small\color{InkGray} 2026 年 9 月}
\end{center}
\vspace{3.0cm}
\end{titlepage}

\newpage

% ---------------------------------------------------------------- 摘要与目录
\pagenumbering{Roman}
% ---------------------------------------------------------------- 摘要（第一部分分册）
% 与 tex/abstract.tex 同源，仅供“第一部分”分册使用：
%   * 保留全报告定位一句与“第一部分的主要发现”四条；
%   * 删去“第二部分的主要发现”四条（其口径、样本与结论属于第二部分，见另册），
%     并去掉指向第二部分章节的两个 \ref（分册中该标签不存在，会印成 ??）。
\section*{摘要}
\addcontentsline{toc}{section}{摘要}

\begin{keybox}[报告要点]
\small
本报告以\textbf{申万行业分类（2021 版）}为统一框架，分为两大部分：
第一部分从供需关系出发分析农业各子行业的周期性，以\hl{猪周期为基准}
测算跨品种、跨行业的相关性与传导时滞；第二部分系统梳理
申万农林牧渔行业\textbf{全部 \num{104} 家 A 股上市公司}近三年的经营与融资画像。
\textbf{本册为第一部分}，对应第一章至第六章；第二部分（逐公司画像与融资需求）
为另一册。

\textbf{第一部分的主要发现}：
（1）本轮农业下行（2024---2025）源于\hl{供给恢复而非需求萎缩}——
玉米连续增产、谷物进口大幅收缩，生猪出栏在产能未出清时继续增长，
白羽肉鸡祖代更新量创历史新高；
（2）品种周期长度由\hl{产能调整的生物学时间}决定，从蛋鸡的 \num{6}---\num{12} 个月
到肉牛的 \num{24}---\num{44} 个月以上，谱系跨越 \num{7} 倍；
（3）在统一口径（月度收益率、同一时间窗口）下，
\hl{跨品种相关性的绝对值普遍低于 \num{0.24}}——花生 \num{+0.48}、粳米 \num{+0.29}、
豆一 \num{+0.26} 三个品种略高于 5\% 显著性阈值，玉米 \num{+0.09}、豆粕 \num{+0.13}
等饲料原料品种均不显著，“猪价带动一切农产品”并非统计事实；
（4）截至 2026 年 9 月，生猪自繁自养已连续六个月处于亏损区间
（猪粮比价月度均值 \num{3.96}---\num{4.47}）、能繁母猪存栏同比转负，
而肉牛价格、蛋鸡存栏与水产投苗已率先反转，\hl{农业周期正处于
“猪链磨底、禽链与牛链上行”的分化阶段}。
\end{keybox}

\vspace{6pt}
\noindent\textbf{关键词}\quad 农业周期；猪周期；跨品种相关性；申万行业分类
\clearpage

\newpage
% 目录页首不再单独排版“目录”标题（页眉已标明），只保留页眉标记
\makeatletter
\renewcommand\tableofcontents{%
  \@mkboth{\contentsname}{\contentsname}%
  \@starttoc{toc}%
}
\makeatother
\tableofcontents
\newpage

% ---------------------------------------------------------------- 图表目录
\phantomsection
\addcontentsline{toc}{section}{图目录}
\listoffigures
\clearpage
\phantomsection
\addcontentsline{toc}{section}{表目录}
\listoftables
\newpage

% ---------------------------------------------------------------- 正文
\pagenumbering{arabic}
\pagestyle{fancy}
\fancyhf{}
\fancyhead[L]{\small\sffamily 中国农业周期性分析报告（第一部分）}
\fancyhead[R]{\small\sffamily\leftmark}
\fancyfoot[C]{\small\thepage}
\renewcommand{\headrulewidth}{0.4pt}
\renewcommand{\headrule}{\hbox to\headwidth{\leaders\hrule height \headrulewidth\hfill}}

\part{农业子行业周期性分析}
\section{研究框架、行业分类与数据基础}
\label{sec:framework}

\subsection{研究背景与目标}
\label{sec:fw-background}

农业是典型的\hl{供给滞后型周期性行业}：生产决策依赖于上一期的价格信号，而生物生长周期、
自然条件与政策干预使供给调整存在刚性时滞，价格因此呈现周期性波动。这一特征在生猪上表现
得最为典型（“猪周期”），但并非生猪独有——糖料、棉花、苹果、红枣、天然橡胶等经济作物同样
存在 3$\sim$7 年的产能周期，蛋鸡、肉禽、水产则存在更短的 1$\sim$3 年周期。

本报告要回答两个层面的问题：

\begin{keybox}[报告的两个核心问题]
\normalsize
\textbf{问题一（周期与相关性）}\quad 中国农业各子行业各自遵循什么样的周期？这些周期之间的
供需联系与价格传导关系如何？以猪周期为基准，其他养殖品种（肉牛、肉羊、禽、水产）与经济
作物是同步、领先还是滞后？\par\vspace{4pt}
\textbf{问题二（上市公司画像）}\quad 近三年农业上市公司的经营与财务状况如何？哪些公司完成
了战略转型？历史融资与未来潜在融资需求有多大？
\end{keybox}

选择以\hl{猪周期为基准}的理由有三：其一，生猪是农业中产值最高、价格波动最剧烈的单一品种，
猪肉产量约占肉类总产量的六成，其价格波动直接决定 CPI 食品项的方向；其二，生猪同时是
\textbf{饲料端传导链条最长}的品种——玉米、豆粕需求的三成以上来自生猪养殖，因此猪周期是
整个农业产业链的需求锚；其三，生猪期货 2021 年 1 月上市后，中国农业终于拥有了一个连续、
高频、可比的生猪价格序列，使得跨品种的定量周期比较成为可能。

\subsection{行业分类口径：申万行业分类（2021 版）}
\label{sec:fw-classification}

本报告统一采用\hl{申万行业分类（2021 版）}。农林牧渔为一级行业（指数代码 801010，共 104 家
A 股上市公司），下设 8 个二级行业、17 个三级行业。选择申万分类而非证监会门类或东财行业，
是因为申万分类对农业内部的养殖与加工链做了更细的切分（例如把“生猪养殖”与“肉鸡养殖”分列），
且自 2021 年起配备了连续的行业指数，便于把\emph{产业周期}与\emph{资本市场表现}对齐研究。

\begin{figure}[htbp]
\centering
\begin{tikzpicture}[
  font=\footnotesize,
  lv1/.style={rectangle, rounded corners=2pt, draw=AgriGreen, line width=1pt,
              fill=AgriGreen, text=white, font=\small\heiti, inner sep=5pt, align=center},
  lv2/.style={rectangle, rounded corners=1.5pt, draw=AgriGreen, line width=0.7pt,
              fill=AgriPale, text=InkGray, font=\footnotesize\heiti, inner sep=3.6pt,
              align=center, text width=1.85cm},
  lv3/.style={rectangle, rounded corners=1pt, draw=AgriLight, line width=0.5pt,
              fill=white, text=InkGray!90, font=\scriptsize, inner sep=2.4pt,
              align=left, text width=1.78cm},
  ed/.style={draw=AgriLight, line width=0.5pt, rounded corners=1pt},
]
\node[lv1] (root) at (0,0) {农林牧渔\\\scriptsize 801010};
\foreach \i/\nm/\cd in {1/种植业/801016, 2/养殖业/801017, 3/饲料/801014,
                        4/农产品加工/801012, 5/动物保健/801018, 6/渔业/801015,
                        7/林业/801011, 8/农业综合/801019}{
  \node[lv2] (L2-\i) at ({(\i-4.5)*2.05}, -1.85) {\nm\\\scriptsize\cd};
  \draw[ed, AgriGreen!70] (root) -- (L2-\i);
}
\node[lv3, anchor=north] (a1) at ([yshift=-0.95cm]L2-1.north) {种子 850111\\粮食种植 850112\\其他种植业 850113\\食用菌 850114};
\node[lv3, anchor=north] (a2) at ([yshift=-0.95cm]L2-2.north) {生猪养殖 850172\\肉鸡养殖 850173\\其他养殖 850174};
\node[lv3, anchor=north] (a3) at ([yshift=-0.95cm]L2-3.north) {畜禽饲料 850142\\水产饲料 850143\\宠物食品 850144};
\node[lv3, anchor=north] (a4) at ([yshift=-0.95cm]L2-4.north) {果蔬加工 850151\\粮油加工 850152\\其他农产品加工 850154};
\node[lv3, anchor=north] (a5) at ([yshift=-0.95cm]L2-5.north) {动物保健 850181};
\node[lv3, anchor=north] (a6) at ([yshift=-0.95cm]L2-6.north) {海洋捕捞 850121\\水产养殖 850122};
\node[lv3, anchor=north] (a7) at ([yshift=-0.95cm]L2-7.north) {林业 850131};
\node[lv3, anchor=north] (a8) at ([yshift=-0.95cm]L2-8.north) {农业综合\\（综合类）};
\foreach \i in {1,...,8}{\draw[ed, AgriLight] (L2-\i) -- (a\i);}
\node[draw=AgriGold, dashed, rounded corners=2pt, fill=AgriGold!8,
      text width=13.8cm, align=left, font=\scriptsize, inner sep=5pt]
  (map) at (0,-6.1)
  {\heiti 第一部分品种映射\quad\rmfamily
   种植业 $\to$ 玉米、强麦、稻米\qquad
   养殖业 $\to$ 生猪、鸡蛋（兼论肉禽、牛羊、水产）\qquad
   饲料 $\to$ 豆粕、菜粕、玉米\\
   农产品加工 $\to$ 豆油、菜油、棕榈油、白糖、棉纱\qquad
   渔业 $\to$ 水产料与鱼粉（无期货，用现货与行业指数）\\
   林业 $\to$ 天然橡胶、纸浆\qquad
   另列林果与油料：苹果、红枣、花生（归入种植业与农产品加工）};
\draw[-{Stealth[length=4pt]}, AgriGold, dashed] (a4.south) to[out=-90,in=90] (map.north);
\end{tikzpicture}
\caption{申万行业分类（2021 版）农林牧渔行业树与本报告品种映射}
\label{fig:sw-tree}
\srcfile{申万宏源研究《申万行业分类标准（2021 版）》；行业代码与成份股数量取自 akshare \texttt{sw\_index\_second\_info} 与 \texttt{sw\_index\_third\_info}。}
\end{figure}

\subsection{周期分析框架：供需、产能与价格}
\label{sec:fw-framework}

\subsubsection{蛛网模型：农业周期的理论内核}

农产品市场的供给对价格的反应存在时滞：\emph{本期的产量由上一期的价格决定}。设第 $t$ 期的
供给与需求分别为
\begin{equation}
Q_t^{s} = a + b\,P_{t-1}, \qquad Q_t^{d} = c - d\,P_t, \qquad b, d > 0,
\label{eq:cobweb}
\end{equation}
市场出清 $Q_t^{s}=Q_t^{d}$ 得价格递推式
\begin{equation}
P_t = \frac{c-a}{d} - \frac{b}{d}\,P_{t-1}.
\end{equation}
其收敛条件为 $\left| -b/d \right| < 1$，即\hl{供给价格弹性小于需求价格弹性时价格收敛}。
农产品恰恰相反：产能一旦建成（母猪、果树、胶园、糖厂压榨能力）难以退出，短期供给弹性极小，
而食品消费的收入弹性低、需求价格弹性同样很小，于是 $b/d$ 接近甚至大于 1，价格以\emph{等幅
或发散振荡}的形式在盈亏线两侧摆动——这正是周期存在的数学根源。

\begin{figure}[htbp]
\centering
\begin{tikzpicture}[font=\footnotesize]
\begin{axis}[
  name=cob,
  width=6.0cm, height=4.9cm,
  xmin=0, xmax=10, ymin=0, ymax=10,
  xtick=\empty, ytick=\empty, axis lines=middle,
  axis line style={InkGray, -{Stealth[length=4pt]}},
  label style={font=\scriptsize},
  clip=false,
]
  \addplot[OilBlue, thick, domain=0.5:9.5]{9.5-0.85*x};
  \node[OilBlue, font=\scriptsize, anchor=west] at (axis cs:8.5,2.6) {$Q^d$};
  \addplot[SoftRed, thick, domain=0.5:9.0]{1.0+0.85*x};
  \node[SoftRed, font=\scriptsize, anchor=west] at (axis cs:8.3,8.4) {$Q^s(P_{t-1})$};
  \draw[InkGray, thick, -{Stealth[length=3pt]}] (axis cs:5.0,5.0) -- (axis cs:5.0,5.25);
  \draw[InkGray, thick, -{Stealth[length=3pt]}] (axis cs:5.0,5.25) -- (axis cs:4.4,5.25);
  \draw[InkGray, thick, -{Stealth[length=3pt]}] (axis cs:4.4,5.25) -- (axis cs:4.4,5.76);
  \draw[InkGray, thick, -{Stealth[length=3pt]}] (axis cs:4.4,5.76) -- (axis cs:5.7,5.76);
  \draw[InkGray, thick, -{Stealth[length=3pt]}] (axis cs:5.7,5.76) -- (axis cs:5.7,4.65);
  \draw[InkGray, thick, -{Stealth[length=3pt]}] (axis cs:5.7,4.65) -- (axis cs:3.5,4.65);
  \draw[InkGray, thick, -{Stealth[length=3pt]}] (axis cs:3.5,4.65) -- (axis cs:3.5,6.5);
  \draw[InkGray, thick, -{Stealth[length=3pt]}] (axis cs:3.5,6.5) -- (axis cs:7.2,6.5);
  \draw[InkGray, thick, -{Stealth[length=3pt]}] (axis cs:7.2,6.5) -- (axis cs:7.2,3.35);
  \node[InkGray, font=\scriptsize, anchor=south east, rotate=90]
    at (axis cs:-0.45,9.9) {价格 $P$};
  \node[InkGray, font=\scriptsize, anchor=west] at (axis cs:9.9,-0.55) {数量 $Q$};
  \node[font=\scriptsize, InkGray, anchor=north east, align=right]
    at (axis cs:9.6,7.2) {蛛网路径：\\等幅振荡};
\end{axis}
\begin{axis}[
  at={(cob.right of south east)}, anchor=left of south west, xshift=14pt,
  width=6.4cm, height=4.9cm, xlabel={时间 $t$}, ylabel={价格 $P_t$},
  xmin=0, xmax=12, ymin=0, ymax=10,
  xtick=\empty, ytick=\empty, axis lines=left,
  axis line style={InkGray},
  label style={font=\scriptsize},
]
  \addplot[AgriGreen, thick, domain=0:12, samples=150]{5+3.4*cos(deg(0.52*x))};
  \node[font=\scriptsize, anchor=west, AgriGreen] at (axis cs:0.2,9.4) {$|b/d|<1$ 收敛};
  \addplot[SoftRed, thick, domain=0:11.4, samples=250]{5+2.8*cos(deg(0.52*x))*1.045^x};
  \node[font=\scriptsize, anchor=west, SoftRed] at (axis cs:5.6,8.9) {$|b/d|\geq1$ 等幅/发散};
  \addplot[dashed, InkGray] coordinates {(0,5) (12,5)};
  \node[font=\scriptsize, anchor=west, InkGray] at (axis cs:9.0,4.4) {长期均衡};
\end{axis}
\end{tikzpicture}
\caption{蛛网模型：供给时滞下的价格振荡（左）与振荡幅度演化（右）}
\label{fig:cobweb}
\srcfile{作者根据式 \eqref{eq:cobweb} 自绘；参数 $b/d$ 分别取 0.7 与 1.045。}
\end{figure}

\subsubsection{四阶段周期划分与产能时钟}

本报告统一用\hl{“产能—价格”双指标}刻画周期阶段，以缓解单一价格指标滞后的问题：价格决定
当期盈利，产能决定未来 6$\sim$18 个月的供给，二者组合即可定位周期位置。

\begin{figure}[htbp]
\centering
\begin{tikzpicture}[font=\footnotesize]
  \draw[AgriLight, line width=0.9pt] (0,0) circle (2.45cm);
  \draw[AgriLight!70, line width=0.5pt, dashed] (-2.45,0) -- (2.45,0);
  \draw[AgriLight!70, line width=0.5pt, dashed] (0,-2.45) -- (0,2.45);
  \foreach \a/\nm/\col in {45/产能去化/AgriGreen, 135/产能底部/AgriGold,
                           225/产能扩张/HogPink, 315/产能顶部/SoftRed}{
    \pgfmathsetmacro\rx{1.55*cos(\a)}
    \pgfmathsetmacro\ry{1.55*sin(\a)}
    \node[circle, draw=\col, fill=\col!12, inner sep=2.8pt, font=\scriptsize\heiti,
          text=\col!75!black] at (\rx,\ry) {\nm};
  }
  \node[font=\scriptsize, align=center] at (0,0) {产能\\时钟};
  \node[font=\scriptsize, align=left, anchor=west] at (3.0,0.9)
    {\heiti 阶段判据\\[3pt]
     \textcolor{AgriGreen}{Q1 产能去化}：价格 $<$ 完全成本，产能指标同比为负\\
     \textcolor{AgriGold}{Q2 产能底部}：价格触底，产能降幅收窄，亏损收窄\\
     \textcolor{HogPink}{Q3 产能扩张}：价格 $>$ 成本，补栏/扩产启动\\
     \textcolor{SoftRed}{Q4 产能顶部}：价格见顶，供给同比高增\\[3pt]
     \textit{产能代理指标}：能繁母猪存栏（生猪）、在产蛋鸡存栏（鸡蛋）、\\
     \textit{植棉/糖料种植面积、胶园开割率、冷库库存（苹果/红枣）}};
\end{tikzpicture}
\caption{以“产能—价格”构建的四阶段周期时钟（猪周期基准模板）}
\label{fig:cycle-clock}
\srcfile{作者根据蛛网模型与各品种产业逻辑自绘。}
\end{figure}

\subsubsection{需求结构与传导链条}

本报告把农产品需求拆为三条链：\textbf{①居民直接消费}（口粮、果蔬与肉类蛋白，受人口、收入
与消费习惯驱动，收入弹性低）；\textbf{②工业原料}（玉米$\to$淀粉与酒精、甘蔗甜菜$\to$制糖、
棉花$\to$纺服、橡胶$\to$轮胎、油脂$\to$生柴与食品加工）；\textbf{③饲料原料}（玉米、豆粕、
菜粕、鱼粉$\to$生猪、禽、水产养殖）。三条链的价格传导速度不同，是理解跨品种周期同步性的
关键。

\begin{figure}[htbp]
\centering
\begin{tikzpicture}[
  font=\scriptsize,
  box/.style={rectangle, rounded corners=1.5pt, draw=InkGray!55, line width=0.6pt,
              fill=LightGray, align=center, inner sep=3.4pt, text width=2.0cm},
  core/.style={rectangle, rounded corners=2pt, draw=HogPink, line width=1.1pt,
               fill=HogPink!12, align=center, inner sep=4pt, text width=2.15cm,
               font=\footnotesize\heiti},
  ar/.style={-{Stealth[length=4pt]}, line width=0.8pt, InkGray},
]
\node[box, fill=OilBlue!10, draw=OilBlue] (cons) at (-5.3,2.4) {居民消费\\收入/人口/替代品};
\node[box, fill=AgriGold!12, draw=AgriGold] (ind) at (0,3.5) {工业原料\\淀粉/制糖/纺服/轮胎};
\node[box, fill=AgriGreen!12, draw=AgriGreen] (feed) at (5.3,2.4) {饲料原料\\玉米/豆粕/菜粕/鱼粉};
\node[core] (hog) at (0,0.3) {生猪养殖\\\textbf{需求锚}};
\node[box, fill=HogPink!8] (poultry) at (4.7,-1.3) {禽：肉鸡/蛋鸡};
\node[box, fill=AquaTeal!12, draw=AquaTeal] (aqua) at (5.0,-3.4) {水产养殖};
\node[box] (ruminant) at (-4.7,-1.3) {肉牛/肉羊};
\node[box] (crop) at (-5.0,-3.4) {种植：玉米/大豆};
\draw[ar] (cons) -- (hog);
\draw[ar] (feed) -- (hog);
\draw[ar] (feed) -- (poultry);
\draw[ar] (poultry) -- (aqua);
\draw[ar] (ind) -- (crop);
\draw[ar, SoftRed, line width=1pt] (hog) -- node[above, font=\tiny, SoftRed] {\ 补栏} (poultry);
\draw[ar] (hog) -- (ruminant);
\draw[ar] (ruminant) -- (crop);
\draw[ar, SoftRed, line width=1pt] (hog) -- (aqua);
\node[draw=SoftRed, dashed, rounded corners=1pt, fill=SoftRed!6, font=\tiny,
      align=left, inner sep=3.5pt, text width=8.6cm] at (0,-4.9)
  {\heiti 关键传导（时滞由第 \ref{sec:corr} 章领先滞后检验校准）\\[1.5pt]
   牛市：猪价$\uparrow$ $\to$ 补栏 $\to$ 玉米/豆粕需求$\uparrow$ $\to$ 饲料涨价 $\to$ 养殖成本$\uparrow$（6$\sim$12 个月）\\
   熊市：猪价$\downarrow$ $\to$ 禽、水产与牛羊肉的替代性需求$\uparrow$（3$\sim$6 个月），同时饲料成本下行改善禽类利润};
\end{tikzpicture}
\caption{农业需求三条链与跨品种周期传导路径}
\label{fig:demand-chain}
\srcfile{作者根据产业链逻辑自绘；箭头粗细表示传导强度，红色箭头为猪周期主导的传导。}
\end{figure}

\subsection{数据来源与处理方法}
\label{sec:fw-data}

\subsubsection{数据来源}

本报告全部数据来自公开渠道，通过 \texttt{akshare} 开源财经数据接口获取并本地留痕。

\begin{table}[htbp]
\centering
\small
\caption{主要数据来源与接口}
\label{tab:data-source}
\begin{tabularx}{\textwidth}{@{}l l X l@{}}
\toprule
\textbf{数据类别} & \textbf{来源} & \textbf{接口} & \textbf{频率} \\
\midrule
期货主力连续行情 & 新浪财经 & \texttt{futures\_main\_sina}（26 个农产品相关品种） & 日 \\
生猪价格指数 & 猪易数据 & \texttt{index\_hog\_spot\_price}（含成交均重） & 日 \\
生猪产业指标 & 大连商品交易所 & \texttt{futures\_hog\_core/cost/supply} & 周/月 \\
期货库存与仓单 & 东方财富、99 期货 & \texttt{futures\_inventory\_em/99} & 周 \\
行业指数与成份股 & 申万宏源 & \texttt{index\_hist\_sw}、\texttt{index\_component\_sw} & 日 \\
公司概况与股本 & 巨潮资讯网 & \texttt{stock\_profile\_cninfo}、\texttt{stock\_share\_change\_cninfo} & 事件 \\
财务报表与指标 & 新浪财经、东方财富 & \texttt{stock\_financial\_abstract}、\texttt{stock\_financial\_analysis\_indicator} & 季 \\
公告全文检索 & 巨潮资讯网 & \texttt{stock\_zh\_a\_disclosure\_report\_cninfo} & 事件 \\
宏观与价格指数 & 国家统计局、中价指数 & \texttt{macro\_china\_cpi\_yearly}、\texttt{macro\_china\_qyspjg} & 月 \\
\bottomrule
\end{tabularx}
\end{table}
\srcfile{akshare 版本 1.18.97；数据抓取脚本见 \texttt{scripts/fetch\_*.py}，抓取时间戳、行列数与接口名记录于 \texttt{data/raw/\_manifest.json}。}

\subsubsection{处理方法}

\begin{enumerate}
  \item \textbf{主力连续价格}的拼接会在换月处产生跳空，本报告对\emph{收益率}类统计
  （波动率、相关系数）使用月度对数收益率，并剔除换月跳空异常值（$|\ln$ 收益$|>35\%$
  视为换月跳空）；对\emph{价格水平}类图形使用月度均价。
  \item \textbf{归一化}：跨品种比较统一以 2021 年 1 月$=100$（生猪期货上市月）为基期；
  数据长度不足的品种以其首个完整月份为基期并在图中注明。
  \item \textbf{相关性}采用皮尔逊相关系数，样本为月度对数收益率；月频 60 个观测下
  5\% 显著性水平对应 $|\rho|\approx0.25$，故低于该阈值的相关性不作因果解读。
  \item \textbf{领先滞后}检验：以生猪期货月收益率与各品种月收益率在 $-12\sim+12$ 个月
  滞后阶上做互相关，取 $|\rho|$ 最大的滞后阶作为“领先 / 同步 / 滞后”判据。
  \item \textbf{财务口径}：报表数据以合并报表为准，2026 年数据为半年度（截至 2026 年
  6 月 30 日）；“近三年”的统计口径为 2023、2024、2025 三个完整会计年度加 2026 年
  半年度。股本变动、分红与公告数据的时间窗口为 2023 年 9 月至 2026 年 9 月。
\end{enumerate}

\section{猪周期：中国农业的核心周期}
\label{sec:hog}

生猪是中国农业中体量最大、价格波动最剧烈、产业链最长的一条主线。它连接着种植（玉米、大豆）、
饲料加工、兽药疫苗、屠宰与肉制品，因此\hl{猪周期不只是生猪养殖自己的周期，而是整个农业需求侧
的节拍器}。本章先解剖猪周期自身的机制与历史，为第 \ref{sec:corr} 章的跨品种比较提供基准。

\subsection{驱动机制：产能—利润—价格三角}
\label{sec:hog-mechanism}

生猪周期之所以稳定存在，是因为生产端存在一个\hl{不可逆且有时滞的产能决定变量——能繁母猪存栏}。
其传导链条为：

\begin{center}
\begin{tikzpicture}[font=\footnotesize,
  b/.style={rectangle, rounded corners=2pt, draw=HogPink, line width=0.9pt,
            fill=HogPink!10, align=center, inner sep=4pt, text width=2.55cm,
            minimum height=1.15cm},
  e/.style={-{Stealth[length=4pt]}, line width=1pt, InkGray}]
  \node[b] (p1) at (0,0) {\textbf{①} 猪价上涨\\养殖利润转正};
  \node[b] (p2) at (3.55,0) {\textbf{②} 补栏：后备母猪\\转为能繁母猪};
  \node[b] (p3) at (7.10,0) {\textbf{③} 约 10 个月后\\生猪出栏增加};
  \node[b] (p4) at (7.10,-2.1) {\textbf{④} 供给过剩\\猪价下跌、亏损};
  \node[b] (p5) at (3.55,-2.1) {\textbf{⑤} 去产能：淘汰\\能繁母猪、清栏};
  \node[b] (p6) at (0,-2.1) {\textbf{⑥} 10$\sim$12 个月后\\供给收缩};
  \draw[e] (p1) -- (p2);
  \draw[e] (p2) -- (p3);
  \draw[e] (p3) -- (p4);
  \draw[e] (p4) -- (p5);
  \draw[e] (p5) -- (p6);
  \draw[e] (p6) -- (p1);
  \node[font=\scriptsize, InkGray, align=center, text width=14.6cm]
    at (3.55,-3.35)
    {关键时滞：\textbf{母猪配种 $\to$ 商品猪出栏约 10 个月}（妊娠 114 天 + 育肥约 6 个月）；
     \textbf{产能调整 $\to$ 价格反应 10$\sim$14 个月}。\\
     由于退出（淘汰母猪）比进入（补栏）耗时更长，周期呈现“涨得快、跌得久”的非对称形态。};
\end{tikzpicture}
\end{center}

在蛛网模型的框架下（第 \ref{sec:fw-framework} 节），生猪的供给价格弹性在半年内接近零（母猪的
生物学周期刚性），需求价格弹性同样很低（猪肉是刚需蛋白），于是 $b/d$ 接近 1，价格呈现
\hl{等幅振荡}：2021 年以来生猪期货月度价格的三次上行段平均持续 5.3 个月、平均涨幅 +41\%，
三次下行段平均持续 15.6 个月、平均跌幅 -46\%——\emph{跌的时间比涨的时间长一倍以上}，
这是产能退出慢于进入的直接结果。

\subsection{四轮完整周期：2015---2026 年}
\label{sec:hog-history}

我们用 ZigZag 算法（波动幅度阈值 20\%）对 2015 年以来的生猪现货价格指数做周期分段，可以清晰
识别出四轮完整周期（谷$\to$谷）。

\begin{figure}[htbp]
\centering
\begin{tikzpicture}
\begin{axis}[agri axis, yearticks,
  width=14.6cm, height=7.0cm,
  xlabel={年份}, ylabel={生猪价格指数},
  xmin=2014.8, xmax=2027.2, ymin=40, ymax=290,
  xtick={2015,2016,2017,2018,2019,2020,2021,2022,2023,2024,2025,2026,2027},
  ytick={50,100,150,200,250},
  legend style={at={(0.02,0.97)}, anchor=north west, cells={anchor=west}},
  legend entries={生猪价格指数（月度）, 周期骨架（ZigZag 20\%）},
]
\addplot[AgriGreen, thick] table[x index=0, y index=1]{data/clean/hog_spot_monthly.dat};
\addplot[SoftRed, very thick, mark=*, mark size=1.5pt]
  table[x index=0, y index=1]{data/clean/hog_zigzag_spot.dat};
% 周期标注
\node[font=\scriptsize, SoftRed, anchor=south] at (axis cs:2016.4,150) {第一轮顶部};
\node[font=\scriptsize, SoftRed, anchor=south] at (axis cs:2020.1,272) {非洲猪瘟};
\node[font=\scriptsize, SoftRed, anchor=west] at (axis cs:2021.75,252) {2021 崩塌};
\node[font=\scriptsize, SoftRed, anchor=south west] at (axis cs:2022.8,186) {2022 反弹};
\node[font=\scriptsize, InkGray, anchor=north] at (axis cs:2026.4,60) {2026 年 4 月新低 63.7};
\end{axis}
\end{tikzpicture}
\caption{生猪现货价格指数与四轮周期（2015 年 1 月---2026 年 9 月，周度转月度）}
\label{fig:hog-price-cycle}
\srcfile{猪易数据生猪价格指数（akshare \texttt{index\_hog\_spot\_price}），作者按月取均值并做 ZigZag 分段。}
\end{figure}

% 由 scripts/make_tables.py 自动生成，请勿手工编辑
\begin{table}[htbp]
\centering
\small
\resizebox{\textwidth}{!}{%
\begin{tabular}{@{}llrrlrr@{}}
\toprule
起点 & 终点 & 起点价 & 终点价 & 方向 & 持续月数 & 幅度\% \\
\midrule
2015-01 & 2015-03 & 93.12 & 84.49 & 下跌 & 1.90 & -9.30 \\
2015-03 & 2016-05 & 84.49 & 143.44 & 上涨 & 14.00 & +69.80 \\
2016-05 & 2018-05 & 143.44 & 70.32 & 下跌 & 24.00 & -51.00 \\
2018-05 & 2020-02 & 70.32 & 266.29 & 上涨 & 21.00 & +278.70 \\
2020-02 & 2020-05 & 266.29 & 207.56 & 下跌 & 3.00 & -22.10 \\
2020-05 & 2020-07 & 207.56 & 259.19 & 上涨 & 2.00 & +24.90 \\
2020-07 & 2020-11 & 259.19 & 204.26 & 下跌 & 4.00 & -21.20 \\
2020-11 & 2021-01 & 204.26 & 248.42 & 上涨 & 2.00 & +21.60 \\
2021-01 & 2021-09 & 248.42 & 88.46 & 下跌 & 8.00 & -64.40 \\
2021-09 & 2021-11 & 88.46 & 120.43 & 上涨 & 2.00 & +36.10 \\
2021-11 & 2022-03 & 120.43 & 84.22 & 下跌 & 4.00 & -30.10 \\
2022-03 & 2022-10 & 84.22 & 183.09 & 上涨 & 7.00 & +117.40 \\
2022-10 & 2023-07 & 183.09 & 98.44 & 下跌 & 9.00 & -46.20 \\
2023-07 & 2024-08 & 98.44 & 141.32 & 上涨 & 13.00 & +43.60 \\
2024-08 & 2026-04 & 141.32 & 63.72 & 下跌 & 19.90 & -54.90 \\
\bottomrule
\end{tabular}
}
\caption{生猪现货价格指数的周期分段（2015 年至今）}
\label{tab:hog-legs-long}
\srcfile{资料来源于 akshare（猪易数据生猪价格指数、新浪财经生猪期货主力连续），分段用 ZigZag 算法（波动幅度阈值 20\%）由作者计算。}
\end{table}


四轮周期的特征可以概括为：

\begin{keybox}[猪周期的四条经验规律]
\normalsize
\textbf{① 上行短、下行长}：上行段平均 8.7 个月、平均涨幅 +\hl{85\%}；下行段平均 9.2 个月、
平均跌幅 -37\%。价格以“快涨慢跌”的方式完成一轮分配。\par\vspace{3pt}
\textbf{② 外生冲击决定振幅}：非洲猪瘟使 2018 年 5 月至 2020 年 2 月的上行幅度达到
+\hl{278.7\%}，是常规周期（+40\%$\sim$+70\%）的四倍以上。\par\vspace{3pt}
\textbf{③ 高利润必然导致产能过冲}：2019---2020 年的超级利润吸引了大量资本进入养殖业，
2021 年 1 月至 9 月价格下跌 -64.4\%，是历史上最陡峭的下跌段。\par\vspace{3pt}
\textbf{④ 本轮下行是历史最长}：2024 年 8 月至 2026 年 4 月连续下行 19.9 个月、跌幅 -54.9\%，
价格指数 63.7 为 2015 年有数据以来最低，超过 2018 年 5 月的 70.3。
\end{keybox}

这一形态与产业界的观察一致：21 世纪以来生猪产业已经历 6 个完整周期，前 5 个周期相对规律、
平均约 48 个月（上行 24 个月、下行 24 个月）；\hl{自第 6 个周期（2022 年）以来周期规律失灵、
下行期显著延长}，第 7 个周期自 2024 年 2 月上行至 2024 年 8 月后进入下行，
至 2025 年底仍未结束\cite{sun2026hog}。本报告的 ZigZag 分段给出的是同一结论的量化版本：
价格波动的\emph{振幅}并未收敛，但\emph{时间结构}已被产能规模化与政策调控改写。

\subsection{产能视角：能繁母猪存栏的四年调整}
\label{sec:hog-capacity}

价格是周期的表象，产能才是周期的本体。中国能繁母猪存栏在 2012---2013 年见顶于 5132 万头，
此后的路径是“阶梯式下移 + 一次断裂”：

\begin{figure}[htbp]
\centering
\begin{tikzpicture}
\begin{axis}[agri axis, yearticks,
  width=14.4cm, height=6.8cm,
  xlabel={年份}, ylabel={能繁母猪存栏（万头）},
  xmin=2008.5, xmax=2026.2, ymin=2800, ymax=5400,
  xtick={2009,2011,2013,2015,2017,2019,2021,2023,2025},
  ytick={3000,3500,4000,4500,5000},
  legend style={at={(0.98,0.97)}, anchor=north east},
  legend entries={年末存栏（年度）, 2025 年季度/月度},
]
\addplot[AgriGreen, thick, mark=square*, mark size=1.6pt]
  table[x index=0, y index=1]{data/clean/hog_能繁母猪存栏_y.dat};
\addplot[SoftRed, very thick, mark=*, mark size=1.8pt, dashed]
  table[x index=0, y index=1]{data/clean/hog_能繁母猪存栏_q.dat};
\node[font=\scriptsize, InkGray, anchor=north east, align=right]
  at (axis cs:2014.2,3450) {产能持续下滑\\（2015---2017）};
\node[font=\scriptsize, SoftRed, anchor=south, align=center]
  at (axis cs:2018.6,3120) {非洲猪瘟\\-24.5\%};
\node[font=\scriptsize, AgriGreen, anchor=north, align=center]
  at (axis cs:2023.2,3980) {产能再度回落\\（2023---2025）};
% 政策目标线（3900 万头）：若能核实则由研究资料补充
\end{axis}
\end{tikzpicture}
\caption{全国能繁母猪存栏：年度数据（2009---2024）与 2025 年季度/月度数据}
\label{fig:hog-capacity}
\srcfile{国家统计局；2025 年数据为季末与月末口径（akshare \texttt{futures\_hog\_supply}）。}
\end{figure}

% 由 scripts/make_tables.py 自动生成，请勿手工编辑
\begin{table}[htbp]
\centering
\footnotesize
\resizebox{\textwidth}{!}{%
\begin{tabular}{@{}lrrrr@{}}
\toprule
期间 & 能繁母猪（万头） & 猪肉产量（万吨） & 生猪存栏（万头） & 生猪出栏（万头） \\
\midrule
2025年一季度（末） & 4\,039 & 1\,602 & 41\,731 & 19\,476 \\
2025年二季度（末） & 4\,043 & 1\,418 & 42\,447 & 17\,143 \\
2025年7月 & 4\,042 & 0 & 0 & 0 \\
2025年8月 & 4\,038 & 0 & 0 & 0 \\
2025年三季度（末） & 4\,035 & 1\,348 & 43\,680 & 16\,373 \\
2025年10月 & 3\,990 & 0 & 0 & 0 \\
\bottomrule
\end{tabular}
}
\caption{2025 年能繁母猪存栏与生猪产销（季度与月度口径）}
\label{tab:hog-capacity-recent}
\srcfile{国家统计局数据（经 akshare \texttt{futures\_hog\_supply} 转载）；0 表示该口径当期未公布。}
\end{table}


三个关键事实：\hl{① 2018 年存栏从 4226 万头降至 3189 万头（-24.5\%）}，2019 年进一步降至
3080 万头，这是非洲猪瘟造成的产能断裂，直接导致 2019---2020 年猪价三倍上涨；
\hl{② 2021---2022 年补栏过冲至 4390 万头}，随后 2023、2024 年连续回落至 4142、4078 万头
（2024 年末为 3900 万头调控目标的 104.6\%，接近绿色区间上限\cite{sun2026hog}）；
\hl{③ 2025 年下半年起进入连续去化阶段}：2025 年末全国能繁母猪存栏 3961 万头，
较上年末减少 117 万头（-2.9\%），2026 年一季度末降至 3904 万头（同比 -3.3\%），
2026 年 7 月进一步调减至 \hl{3780 万头}，重新回到产能调控的绿色合理区间
\cite{sun2026hog,caaa2026q1,jiemian2026pig}。仍然偏高的原因是
\emph{母猪生产效率（PSY）持续提升}：2025 年全国经产母猪年均提供断奶仔猪 23.4 头、
同比增加 2.4 头，头部企业 PSY 超过 29 头\cite{sun2026hog}；同样存栏对应更高的猪肉产出，
猪肉产量从 2021 年 5296 万吨升至 2023 年 5794 万吨（+9.4\%），
2024 年 5706 万吨、2025 年 \hl{5938 万吨（+4.1\%，历史新高）}\cite{sun2026hog}
——这正是本轮猪价长期低迷的根本原因。

\subsection{成本与利润：猪粮比价的信号意义}
\label{sec:hog-cost}

生猪养殖利润是产能调整的唯一信号，而利润由“猪价 $-$ 饲料成本”决定。猪粮比价（生猪出栏价与
玉米价格之比）是官方预警指标，\hl{低于 6:1 进入三级预警、低于 5:1 进入一级预警}。

\begin{figure}[htbp]
\centering
\begin{tikzpicture}
\begin{groupplot}[group style={group size=3 by 1, horizontal sep=0.7cm},
  agri axis, width=4.65cm, height=4.2cm,
  tick label style={font=\tiny}, label style={font=\scriptsize},
  every axis plot/.append style={line width=0.9pt}]
\nextgroupplot[yearticks, title={（a）猪粮比价（倍）},
  xmin=2020, xmax=2027, ymin=0, ymax=20,
  xtick={2020,2022,2024,2026}, ytick={0,5,10,15,20}]
  \addplot[SoftRed, thick] table[x index=0, y index=1]{data/clean/hogx_猪粮比价.dat};
  \addplot[dashed, InkGray, domain=2020:2027]{6};
  \addplot[dashed, SoftRed, domain=2020:2027]{5};
  \node[font=\tiny, InkGray, anchor=west] at (axis cs:2020.1,6.6) {6:1 三级预警};
  \node[font=\tiny, SoftRed, anchor=west] at (axis cs:2020.1,3.4) {5:1 一级预警};
\nextgroupplot[yearticks, title={（b）仔猪价格（元/公斤）},
  xmin=2019, xmax=2027, ymin=0, ymax=160,
  xtick={2019,2021,2023,2025}]
  \addplot[Grain, thick] table[x index=0, y index=1]{data/clean/hogx_仔猪价格.dat};
\nextgroupplot[yearticks, title={（c）二元母猪价格（元/公斤）},
  xmin=2018, xmax=2027, ymin=0, ymax=90,
  xtick={2018,2020,2022,2024,2026}]
  \addplot[HogPink, thick] table[x index=0, y index=1]{data/clean/hogx_二元母猪价格.dat};
\end{groupplot}
\end{tikzpicture}
\caption{生猪养殖成本与利润的关键指标：猪粮比价、仔猪价格与二元母猪价格}
\label{fig:hog-cost}
\srcfile{玄田数据（akshare \texttt{futures\_hog\_supply}、\texttt{futures\_hog\_cost}）；月度均值。}
\end{figure}

三点结论：

\begin{enumerate}
  \item \textbf{猪粮比价已连续六个月处于一级预警区间}。2026 年 3---8 月的月度均值分别为
  4.47、4.03、4.12、\hl{3.96}、4.37、4.44，均在 5:1 以下，6 月的 3.96:1 为本轮最低；
  作为对比，2024、2025 年的年均值分别为 7.15 与 6.27。需要说明的是，该序列自 2020 年
  1 月起有数据，在此之前仅 2022 年 3---4 月与 2023 年 7 月短暂跌破 5:1。行业自 2026 年
  一季度起进入\emph{全面现金亏损}状态。
  \item \textbf{仔猪价格是本轮周期最灵敏的领先指标}。仔猪价格从 2020 年 8 月的峰值
  108.6 元/公斤回落至 2026 年 8 月的 \hl{22.05 元/公斤}，跌幅约 80\%，已回到 2019 年
  1 月非洲猪瘟最严重时期（21.6 元/公斤）的水平；仔猪价格反映的是\emph{补栏意愿}，
  其持续低迷说明养殖户对未来 6 个月后的猪价预期悲观，是产能继续去化的确认信号。
  \item \textbf{二元母猪价格同步走弱}（2026 年 6 月 26.8 元/公斤，为 2018 年该序列有数据
  以来的最低值，2020 年峰值为 75.6 元/公斤），说明母猪更新补栏同样停滞，
  产能去化已从“淘汰落后母猪”进入“停止正常更新”的第二阶段。
\end{enumerate}

\subsection{疫病与政策：两次外生冲击}
\label{sec:hog-shock}

猪价的三倍暴涨（2019---2020）与五年阴跌（2021---2026）分别对应两次外生冲击：
一次是\hl{疫病冲击}，一次是\hl{政策与结构冲击}。

\textbf{（1）疫病：非洲猪瘟（2018 年 8 月起）。}2018 年 8 月沈阳报告首例非洲猪瘟疫情后，
产能以断崖方式退出——能繁母猪存栏从 2017 年的 4226 万头降至 2018 年的 3189 万头、
2019 年的 3080 万头；由于生猪的生物学生产周期不可压缩，供给缺口无法在 12 个月内回补，
价格上行持续了 21 个月、涨幅接近 2.8 倍。疫情同时永久改变了产业组织形态：
全国生猪养殖规模化比重从 2018 年的 49.1\% 升至 2025 年的约 \hl{73.0\%}，
前 10 位企业出栏量占比从 8.1\% 升至 \hl{29.7\%}，散养户数量从 2018 年末的 2706 万户
降至 2025 年末的约 1672 万户（-38.2\%）\cite{sun2026hog}。防疫能力取代成本成为第一竞争力，
这也是 2021 年之后供给恢复速度快于历史任何一轮周期、价格下行时间被拉长的结构性原因。

\textbf{（2）政策：从“预警”到“调控”。}2021 年 6 月多部委联合印发《完善政府猪肉储备
调节机制 做好猪肉市场保供稳价工作预案》，确立以猪粮比价、能繁母猪存栏量变化率、
36 个大中城市精瘦肉零售价为核心的三级预警体系：猪粮比低于 6:1 发布三级预警，
连续 3 周处于 5:1$\sim$6:1（或能繁母猪存栏单月同比降幅达 5\%）发布二级预警，
低于 5:1（或能繁母猪存栏单月同比降幅达 10\%）发布一级预警\cite{ndrc2021reserve,nbd2025pigcorn}。
预案的调控靶心是可量化、可备案的\emph{能繁母猪存栏}：

\begin{itemize}
  \item \textbf{2024 年修订}：全国能繁母猪正常保有量目标由 4100 万头下调至
  \hl{3900 万头}，绿色区域下限由正常保有量的 95\% 调至 92\%
  \cite{moa2024capacity}；
  \item \textbf{2026 年修订}（2026 年 5 月 14 日印发）：目标进一步下调至
  \hl{3750 万头左右}，绿色区域收窄为 92\%$\sim$103\%，上限由 105\% 下调至 103\%，
  对应月度存栏约 3450 万$\sim$3863 万头；同时首次把能繁母猪存栏总量 10 万头以上的
  大型集团企业纳入全国产能调控监测名单、实行年度生产备案管理\cite{stcn2026capacity}；
  \item \textbf{“四个带头”}（2026 年 6 月 17 日，农业农村部与国家发改委联合座谈）：
  对大型养殖企业提出带头压减产能产量、带头严控二次育肥、带头淘汰弱仔猪、
  带头降低出栏体重的刚性要求\cite{jiemian2026pig}。
\end{itemize}

政策调控的思路可概括为“\emph{长期调母猪，中期调仔猪，短期调肥猪}”\cite{stcn2026capacity}。
从效果看，2025 年 7 月以来全国能繁母猪存栏连续下降，2026 年 3 月新生仔猪数量
17 个月后首次同比下降\cite{stcn2026capacity}，政策确实加快了产能去化的节奏——
但\hl{调控的目标不是消灭周期，而是压低周期的振幅}，这一点在下文的跨品种分析中还会反复出现。

\subsection{效率变量：出栏均重与二次育肥}
\label{sec:hog-weight}

同样的存栏可以对应不同的猪肉供给量，\hl{出栏体重是产能之外的第二个供给调节阀}。
当价格预期上行时，养殖户倾向于压栏（二次育肥），体重上升会短期减少出栏、推高价格，
但在 2---3 个月后转化为更大的集中出栏压力；价格预期下行时则加速出栏、体重快速回落。

\begin{figure}[htbp]
\centering
\begin{tikzpicture}
\begin{axis}[agri axis, yearticks,
  width=13.6cm, height=5.6cm,
  xlabel={年份}, ylabel={生猪出栏平均体重（公斤）},
  xmin=2015, xmax=2027, ymin=95, ymax=150,
  xtick={2015,2016,2017,2018,2019,2020,2021,2022,2023,2024,2025,2026,2027},
  ytick={100,110,120,130,140,150},
  legend pos=south west, legend entries={出栏平均体重（月度）},
]
\addplot[AquaTeal, thick] table[x index=0, y index=1]{data/clean/hog_spot_均重.dat};
\addplot[dashed, InkGray, domain=2015:2027]{120};
\node[font=\scriptsize, InkGray, anchor=south west] at (axis cs:2015.1,120.5) {120 公斤};
\node[font=\scriptsize, SoftRed, anchor=south west] at (axis cs:2021.1,146.0) {2021 年压栏峰值 146.2};
\node[font=\scriptsize, SoftRed, anchor=north east, align=right]
  at (axis cs:2026.5,138) {2026 年 6---7 月\\回升至 139.2};
\node[font=\scriptsize, AgriGreen, anchor=west, align=left]
  at (axis cs:2026.55,124) {8---9 月\\骤降至 125.3};
\end{axis}
\end{tikzpicture}
\caption{生猪出栏平均体重（2015 年 8 月---2026 年 9 月）}
\label{fig:hog-weight}
\srcfile{猪易数据成交均重（akshare \texttt{index\_hog\_spot\_price}），月度均值；2015 年早期观测缺失。}
\end{figure}

体重数据的三个要点：\hl{① 2021 年中的 146.2 公斤}是压栏行为的极端值，出现在猪价崩塌前，
随后成为价格加速下跌的助推器；\hl{② 年均值从 2023 年的 109.5 公斤逐级抬升至 2025 年 127.3
公斤、2026 年 127.4 公斤}，说明行业平均出栏体重在周期下行期反而上升（成本压力下以体重摊薄
固定成本）；\hl{③ 2026 年 6---7 月再度回升至 139.2 公斤后，8---9 月骤降至 125.3 公斤}
（两个月下降约 14 公斤），这是年内最剧烈的降重，对应的是养殖户在深度亏损下的集中抛售——
它意味着\emph{短期供给仍在被强制出清}，但同时也在为 2027 年的供给缺口创造条件。

\subsection{当前周期定位}
\label{sec:hog-position}

综合产能、价格、成本三个维度，2026 年 9 月的生猪产业处于\hl{四阶段时钟的 Q1（产能去化）末段、
接近 Q2（产能底部）}：

\begin{itemize}
  \item 价格维度：价格指数 63.7（2026 年 4 月）为本轮最低，2026 年 9 月小幅回升至 75.8
  （月度均值）；2026 年 4 月全国活猪均价 10.07 元/公斤，
  2026 年 7 月回升至 11.50 元/公斤，仍低于 2025 年全国平均养殖成本
  14.1 元/公斤\cite{dxzq2026pig,jiemian2026pig,sun2026hog}；
  \item 产能维度：能繁母猪存栏从 2022 年的 4390 万头降至 2026 年 7 月的
  \hl{3780 万头}（-13.9\%），已回到 3750 万头新调控目标的绿色区间\cite{jiemian2026pig}；
  \item 效率维度：2025 年样本出栏均重约 127.13 公斤、为近几年最高，
  2026 年 8---9 月骤降至 125.3 公斤，说明政策要求下的降重已经落地
  \cite{stcn2026profit}；
  \item 行为维度：仔猪价格与二元母猪价格处于历史低位，补栏意愿尚未恢复。
\end{itemize}

历史规律是：\emph{价格在产能去化完成前 6---9 个月见底}。本轮的特殊性在于三点：
一是产能去化由\hl{政策推动而非亏损自发出清}（2026 年修订方案的目标下调与“四个带头”），
去化速度比历史周期更快；二是能繁母猪存栏的绝对水平已回到绿色区间，
但 PSY 提升与规模化企业的成本优势使\emph{有效供给}仍高于名义存栏所隐含的水平；
三是需求端偏弱——2025 年居民家庭人均猪肉消费量 26.6 公斤、同比下降 5.4\%，
连续两年下跌\cite{sun2026hog}。因此 2026 年下半年的判断是：
\hl{价格底部大概率已在 2026 年上半年出现，但回升力度受制于需求，
“长期磨底”而非“快速反转”是基准情形}。农业农村部 2026 年 7 月的官方研判亦为
“供需关系逐步改善、但供强需弱格局尚未根本改变”\cite{jiemian2026pig}。

\section{粮食与油料：玉米、大豆及压榨链}
\label{sec:grain}

如果说猪周期是农业需求侧的节拍器，那么\hl{粮食与油料就是被这个节拍器驱动、但有自己的旋律}的
另一条主线。中国农业的“粮食—饲料—养殖”链条中，玉米与大豆是两个枢纽品种：前者是国内唯一
能够横跨饲料与深加工两大需求的大宗粮食品种，后者则以超过 1 亿吨的年进口量把中国农业与
南美、北美的天气和贸易政策直接绑定。

\subsection{产业链结构：两类枢纽品种}
\label{sec:grain-structure}

\begin{center}
\begin{tikzpicture}[font=\footnotesize,
  n/.style={rectangle, rounded corners=2pt, draw=Grain!80!black, line width=0.8pt,
            fill=Grain!12, align=center, inner sep=3.5pt, minimum height=1.05cm},
  e/.style={-{Stealth[length=3.5pt]}, line width=0.9pt, InkGray},
  lab/.style={font=\scriptsize, InkGray, align=center}]
  % 玉米链
  \node[n, text width=1.5cm] (cg) at (0,2.6) {玉米\\30123 万吨};
  \node[n, text width=1.9cm] (cf) at (2.9,2.6) {饲料用\\65\%$\sim$70\%};
  \node[n, text width=1.7cm] (ci) at (5.9,2.6) {深加工\\25\%$\sim$30\%};
  \node[n, text width=1.7cm] (cl) at (8.9,2.6) {生猪养殖\\饲用中约 50\%};
  \draw[e] (cg) -- (cf);
  \draw[e] (cg) -- (ci);
  \draw[e] (cf) -- (cl);
  % 大豆链
  \node[n, text width=1.9cm] (sg) at (0,0.4) {大豆\\进口 11183 + 国产 2091};
  \node[n, text width=1.6cm] (sp) at (3.1,0.4) {压榨\\10150 万吨};
  \node[n, text width=1.5cm] (sm) at (6.2,0.4) {豆粕\\约 80\%};
  \node[n, text width=1.5cm] (so) at (8.6,0.4) {豆油\\约 18\%};
  \node[n, text width=1.7cm] (sj) at (6.2,-1.8) {生猪、家禽\\饲料（刚性需求）};
  \draw[e] (sg) -- (sp);
  \draw[e] (sp) -- (sm);
  \draw[e] (sp) -- (so);
  \draw[e] (sm) -- (sj);
  \node[lab, anchor=west] at (-1.6,3.5) {（a）玉米：两条需求腿};
  \node[lab, anchor=west] at (-1.6,1.3) {（b）大豆：压榨链};
\end{tikzpicture}
\end{center}

\textbf{玉米的两条需求腿。}玉米饲用消费占总消费量的 \hl{65\%$\sim$70\%}，其中生猪养殖消耗
约占饲用消费的 50\%、禽料约 40\%；深加工年耗玉米约 7800 万吨、占 25\%$\sim$30\%，
全国深加工产能已超过 1.25 亿吨\cite{lswz2025corn}。这一结构决定了玉米价格同时受
\emph{养殖产能周期}与\emph{深加工开工率}影响，而两者的周期长度完全不同：
养殖产能是 3$\sim$4 年的长周期，深加工利润则是 1 年以内的短周期。

\textbf{大豆的压榨链。}大豆压榨产出约 80\% 的豆粕与 18\% 的豆油，豆粕几乎全部用于饲料。
2025/26 年度中国大豆压榨量约 10150 万吨\cite{casde123}，压榨产能 18963 万吨、近五年
复合增长率 2.33\%\cite{zhuochuang2026soy}。由于豆粕需求来自生猪与家禽存栏，
\hl{压榨量与猪周期同向}，但豆粕的\emph{价格}由进口成本（CBOT + 升贴水 + 汇率）决定，
这是理解“量与价背离”的关键。

\subsection{供给端：产量、进口与政策}
\label{sec:grain-supply}

\subsubsection{玉米：连续丰产与进口政策转向}

\begin{table}[htbp]
\centering
\small
\begin{tabular}{@{}llr@{}}\toprule
指标 & 最新值 & 备注 \\
\midrule
播种面积 & 44960.8 千公顷 & 2025 年（玉米） \\
产量 & 30123.5 万吨 & 2025 年，单产 6700.0 公斤/公顷 \\
进口量 & 264.77 万吨 & 2025 年日历年，同比 -80.8\% \\
进口量 & 182.4 万吨 & 2024/25 年度（2024 年 10 月---2025 年 9 月），同比 -92\% \\
进口配额 & 720 万吨 & 2004 年至今未变，配额内关税 1\%、配额外 65\% \\
替代品进口 & 2301 万吨 & 2024/25 年度高粱、大麦、DDGS、木薯干合计，同比 -54\% \\
\bottomrule
\end{tabular}
\caption{中国玉米供给端关键指标（2024---2025）}
\label{tab:corn-supply}
\end{table}

玉米供给端最重要的事件是\hl{进口的断崖式退出}：在 2021---2023 年连续三年进口超过
2000 万吨之后，2024 年进口回落至 1364 万吨（-49.7\%），
2024/25 年度仅 182.4 万吨（-92\%，近十四年最低），其中美玉米仅 3.9 万吨（-98.7\%）、
巴西玉米 72.5 万吨（-95.2\%）\cite{xinhu2025corn}。其直接原因是
2024 年 4 月前后进口政策收紧，月度进口量从百万吨级降至 20 万吨以下\cite{xinhu2025corn}。

这一转变的政策含义是：\hl{国内玉米价格的政策底由“进口成本”变成了“国内种植成本 + 补贴”}。
进口配额 20 余年稳定在 720 万吨、配额外 65\% 的关税，意味着只要进口量被行政性压低，
国内玉米价格就失去了外部的天花板与地板，\emph{价格更多由国内产量与库存决定}。
2026/27 年度预测产量 \num{30600} 万吨（+1.6\%）\cite{casde119}；
2025/26 年度饲用消费估计值上调至 \num{21090} 万吨、工业消费下调至 \num{8100} 万吨\cite{casde119}，
供需处于紧平衡状态。

\subsubsection{大豆：创纪录进口与来源结构重构}

2025 年中国大豆进口 \hl{11183 万吨}（同比 +6.47\%，历史新高），
其中巴西 8232 万吨（占 73.61\%，同比提升 2.54 个百分点）、
美国 1682 万吨（占 15.04\%，同比下滑 6.03 个百分点，为 2015 年有数据以来最低）、
阿根廷 789 万吨（+92.4\%）；进口均价 450 美元/吨、同比下降 10.35\%\cite{zhuochuang2026soy,smm2026soy}。
2025 年 12 月自美国进口大豆为零，为连续第四个月零进口\cite{smm2026soy}。
国产大豆 2025 年产量 2091 万吨（+1.3\%）\cite{casde119}，仅相当于进口量的 19\%，
\hl{大豆是农业领域进口依赖度最高、也最容易被贸易政策扰动的品种}。

\begin{keybox}[供给端的三种冲击类型]
\normalsize
\textbf{① 气候冲击（产量）}：玉米、大豆的国内产量由播种面积与单产决定，单产受生长期
（东北 7---8 月）降水与积温影响；大豆、玉米的主产区高度重合，因此\emph{同一场天气会同时影响两个品种}，
但它们对价格的影响并不同步（大豆国内产量占比低，气候对价格的影响远弱于进口政策）。\par\vspace{3pt}
\textbf{② 政策冲击（进口与种植结构）}：玉米进口政策收紧、大豆扩种政策、玉米大豆生产者补贴、
三大主粮完全成本保险与种植收入保险全国覆盖\cite{cccp2024doc1}，
构成“国内价格地板”；而中央一号文件“深入推进粮油作物大面积单产提升行动”
与“种业振兴”\cite{cccp2025doc1}则决定长期的供给斜率。\par\vspace{3pt}
\textbf{③ 外生冲击（贸易与疫情）}：中美贸易摩擦对美豆份额的影响（2025 年美豆占比降至 15\%）、
以及非洲猪瘟这类疫情对\emph{下游需求}的冲击，是价格短期剧烈波动的主要来源。
\end{keybox}

\subsection{需求端：饲料工业的规模与结构}
\label{sec:grain-demand}

粮食与油料的需求侧可以用一个数字概括：\hl{2025 年全国工业饲料总产量 34225.3 万吨，
同比增长 8.6\%，创历史新高}\cite{chinafeed2025}。分品种看，
猪饲料 16639.4 万吨（+15.6\%，占总量的 48.6\%）、肉禽料 10097.7 万吨（+3.5\%，首次突破 1 亿吨）、
蛋禽料 3282.0 万吨（+1.4\%）、水产料 2323.1 万吨（+2.7\%）、反刍料 1475.8 万吨（+1.8\%）、
宠物料 188.4 万吨（+17.9\%）\cite{chinafeed2025}。上述六类合计 34006 万吨、占总量 \num{99.4}\%，
其余为其他饲料。

\begin{figure}[htbp]
\centering
\begin{tikzpicture}
\begin{axis}[agri axis, yearticks,
  width=13.6cm, height=5.8cm,
  xlabel={年份}, ylabel={全国工业饲料产量（万吨）},
  xmin=2022.6, xmax=2026.2, ymin=12000, ymax=35000,
  xtick={2023,2024,2025}, ytick={15000,20000,25000,30000,35000},
  % 图例移到图外上方：两条序列（总产量、猪饲料）数值区间不同，图内四角都会被数据占用
  legend style={at={(0.5,1.02)}, anchor=south, legend columns=2, cells={anchor=west}},
  legend entries={工业饲料总产量, 其中：猪饲料},
]
\addplot[AgriGreen, thick, mark=*, mark size=2pt] coordinates {
  (2023,32162.7) (2024,31503.1) (2025,34225.3)};
\addplot[HogPink, thick, mark=square*, mark size=2pt] coordinates {
  (2023,14975.2) (2024,14391.3) (2025,16639.4)};
\node[font=\scriptsize, InkGray, anchor=south] at (axis cs:2024,31700) {-2.1\%（近十年首降）};
\node[font=\scriptsize, AgriGreen, anchor=south] at (axis cs:2025,34350) {+8.6\%};
\end{axis}
\end{tikzpicture}
\caption{全国工业饲料产量与猪饲料产量（2023---2025）}
\label{fig:feed-output}
\end{figure}

饲料产量的波动本身就是一次“需求侧周期”的完整案例：2024 年受养殖业深度亏损影响，
全国饲料产量出现近十年首次下降（-2.1\%，其中猪料 -3.9\%、反刍料 -13.3\%、水产料 -3.5\%）\cite{chinafeed2024}；
2025 年在能繁母猪与家禽存栏高位的支撑下，猪料大幅回升 15.6\%，带动总量创新高\cite{chinafeed2025}。
需求侧的第二个结构性变化是\hl{豆粕减量替代}：2025 年饲料企业中豆粕在配合饲料和浓缩饲料中的
占比为 13.4\%、与上年持平，菜粕与棉籽粕等杂粕用量增加 3.0\%\cite{chinafeed2025}；
2023 年豆粕用量曾下降 11.8\%、2024 年再降 4.7\%\cite{chinafeed2024}。
这一趋势直接削弱了“生猪存栏$\to$豆粕需求$\to$豆粕价格”的传导强度。

\subsection{价格周期：玉米与豆粕的不同驱动}
\label{sec:grain-cycle}

\begin{figure}[htbp]
\centering
\begin{tikzpicture}
\begin{axis}[agri axis, yearticks,
  width=14.2cm, height=6.8cm,
  xlabel={年份（2021 年 2 月 $=100$）}, ylabel={归一化价格指数},
  xmin=2021.0, xmax=2026.9, ymin=35, ymax=200,
  xtick={2021,2022,2023,2024,2025,2026},
  % 图例移到图外下方：留在图内会遮住曲线末端的数值标注
  legend style={at={(0.5,-0.24)}, anchor=north, legend columns=5,
                cells={anchor=west}, font=\tiny},
  legend entries={生猪, 鸡蛋, 玉米, 豆粕, 豆油, 棕榈油, 白糖, 棉花, 天然橡胶},
]
\addplot[HogPink, very thick] table[x=x, y=LH]{data/clean/fut_norm_grain.dat};
\addplot[AgriGold, thick] table[x=x, y=JD]{data/clean/fut_norm_grain.dat};
\addplot[Grain, thick] table[x=x, y=C]{data/clean/fut_norm_grain.dat};
\addplot[OilBlue, thick] table[x=x, y=M]{data/clean/fut_norm_grain.dat};
\addplot[AgriGreen, thick] table[x=x, y=Y]{data/clean/fut_norm_grain.dat};
\addplot[violet, thick] table[x=x, y=P]{data/clean/fut_norm_grain.dat};
\addplot[SoftRed, thick] table[x=x, y=SR]{data/clean/fut_norm_grain.dat};
\addplot[AquaTeal, thick] table[x=x, y=CF]{data/clean/fut_norm_grain.dat};
\addplot[InkGray, thick] table[x=x, y=RU]{data/clean/fut_norm_grain.dat};
\node[font=\scriptsize, HogPink, anchor=east] at (axis cs:2026.88,196) {生猪 41.8};
\node[font=\scriptsize, Grain, anchor=east] at (axis cs:2026.88,180) {玉米 80.3};
\end{axis}
\end{tikzpicture}
\caption{主要农产品期货价格归一化走势（2021 年 2 月 $=100$，月度均价；图中数值为 2026 年 9 月的最新指数）}
\label{fig:fut-norm}
\end{figure}

图 \ref{fig:fut-norm} 呈现了本轮周期最关键的\hl{分化}：2021 年 2 月至 2026 年 9 月，
生猪期货价格下跌 58.3\%（指数 41.8），而玉米仅下跌 19.7\%、豆粕下跌 2.5\%、
豆油上涨 11.6\%、棕榈油上涨 43.8\%、天然橡胶上涨 24.8\%。
\emph{下游养殖品价格腰斩，而上游原料价格基本持平甚至上涨}——这直接说明
猪价下跌并未通过饲料需求传导为原料价格下跌，原因有三：一是生猪存栏的绝对水平并未崩塌
（2025 年生猪出栏 71973 万头、同比 +2.4\%），饲料需求仍在增长；
二是原料价格由国际供给与政策主导；三是豆粕减量替代降低了单位生猪的豆粕消耗。

\begin{figure}[htbp]
\centering
\begin{tikzpicture}
\begin{axis}[agri axis, catbar,
  width=13.0cm, height=7.4cm,
  xlabel={年化波动率（\%）}, xmin=0, xmax=30,
  yticklabels from table={data/clean/vol_bar.dat}{label},
  ytick=data,
  legend style={at={(0.97,0.03)}, anchor=south east},
  legend entries={主窗口 2021-02\textasciitilde2026-08, 全样本},
]
% 两组条形错开 0.2 行：同 y 绘制会互相遮盖，只剩较长的一根可见
\addplot[xbar, bar width=3.6pt, fill=AgriGreen!75, draw=AgriGreen]
  table[x=vol, y expr=\coordindex-0.2]{data/clean/vol_bar.dat};
\addplot[xbar, bar width=3.6pt, fill=AgriGold!80, draw=AgriGold]
  table[x=full, y expr=\coordindex+0.2]{data/clean/vol_bar.dat};
\end{axis}
\end{tikzpicture}
\caption{农产品期货年化波动率：主窗口与全样本对比（主窗口为 2021 年 2 月---2026 年 8 月的 67 个月度点；全样本自各品种上市首月至 2026 年 8 月；月度对数收益率年化，$\times\sqrt{12}$）}
\label{fig:vol-bar}
\end{figure}

波动率的分层同样值得注意（图 \ref{fig:vol-bar}）：生猪（27.1\%）、红枣（25.9\%）、
棕榈油（23.7\%）、鸡蛋（22.9\%）构成本轮周期的高波动组，
而粳米（4.0\%）、白糖（8.7\%）、玉米（9.2\%）波动率最低。
\hl{玉米的低波动是“政策定价”的直接结果}——进口受配额约束、国储与地方储备调节库存，
价格被限制在种植成本与替代品价格之间；而生猪、鸡蛋、红枣完全由市场供需定价，
且供给调整存在刚性，因此其波动率约为玉米的三倍（生猪 27.1\% / 玉米 9.2\%）。

\subsection{与猪周期的连接：量与价的分离}
\label{sec:grain-hog-link}

把上述供给与需求拼起来，可以得到本章的核心判断：

\begin{keybox}[粮食油料与猪周期的连接方式]
\normalsize
\textbf{（1）需求侧同向、价格侧弱相关。}饲料产量（尤其是猪料）与生猪存栏高度同向，
2025 年猪料增长 15.6\% 就是生猪存栏高位的直接结果；但玉米、豆粕的\emph{价格}
与生猪价格的相关性很弱（第 \ref{sec:corr} 章的月度收益率相关系数分别为
玉米 $+0.09$、豆粕 $+0.13$，均未达到 $|\rho|>0.24$ 的 5\% 显著性阈值），
因为原料价格由国际供给与政策决定。\par\vspace{3pt}
\textbf{（2）成本传导是非对称的。}原料价格上涨会立即压缩养殖利润（玉米与豆粕占饲料成本
超过 70\%）\cite{sun2026hog}，但养殖亏损压缩产能后，原料需求的下降幅度有限
（存栏的下行慢于利润的下行）。这构成养殖业“利润周期”比“价格周期”更剧烈的原因。\par\vspace{3pt}
\textbf{（3）政策是两条链的共同外生变量。}同一份中央一号文件既规定粮食产能目标
（“确保粮食产量保持在 1.3 万亿斤以上”）\cite{cccp2024doc1}，又规定生猪产能调控机制，
使粮食与养殖在\emph{数量}上被同时锁定、在\emph{价格}上被各自隔离。
\end{keybox}

\section{软商品与经济作物：白糖、棉花、苹果、红枣、花生、天然橡胶}
\label{sec:soft}

本章的七个品种（白糖、棉花、苹果、红枣、花生、天然橡胶、纸浆）与第 \ref{sec:grain} 章的
粮食油料有本质区别：\hl{它们的周期几乎不由中国的需求决定，而由产地气候、全球供需与
政策制度决定}。这决定了它们的周期性形态——不是猪周期那样的“产能$\to$价格”闭环，
而是\emph{外生冲击主导的脉冲式周期}。

\subsection{三类周期来源}
\label{sec:soft-drivers}

\begin{center}
\begin{tikzpicture}[font=\footnotesize,
  b/.style={rectangle, rounded corners=3pt, draw=AgriGreen, line width=0.9pt,
            fill=AgriPale, align=left, inner sep=5pt, text width=4.25cm,
            minimum height=2.6cm},
  h/.style={font=\small\heiti, AgriGreen, anchor=west}]
  \node[b] (t1) at (0,0)
    {\textbf{① 政策定价型}\\[2pt] 白糖、棉花\\[2pt]
     \scriptsize 目标价格补贴、进口配额/滑准税、国储与地方储备共同构成价格区间；
     周期表现为“政策调整$\to$种植面积变化$\to$产量变化”的 3$\sim$4 年节奏，
     价格振幅被制度性压缩。};
  \node[b, right=0.5cm of t1] (t2)
    {\textbf{② 产地气候型}\\[2pt] 苹果、红枣、花生\\[2pt]
     \scriptsize 供给高度集中于单一产区（新疆、陕西、山东），
     花期霜冻、成熟期连阴雨与“大小年”直接决定产量；
     周期表现为“减产$\to$价格暴涨$\to$补种/大年$\to$价格回落”的脉冲。};
  \node[b, right=0.5cm of t2] (t3)
    {\textbf{③ 全球供需型}\\[2pt] 天然橡胶、纸浆\\[2pt]
     \scriptsize 供给在东南亚/全球浆厂、需求在轮胎与造纸；
     中国的角色是消费方，价格由全球产能周期（胶树 7$\sim$10 年生长周期、浆厂投产节奏）
     与国际报价决定。};
  \node[h] at (-1.9,1.7) {周期来源三分法};
\end{tikzpicture}
\end{center}

\subsection{白糖：全球过剩与短缺的切换}
\label{sec:soft-sugar}

白糖的国内供给可以精确到榨季：2023/24 榨季产糖 996 万吨，
2024/25 榨季 \hl{1116.21 万吨（同比 +12.03\%）}，实现连续三个榨季的恢复性增长\cite{sugarassoc2025}；
2025/26 榨季最终统计产量 \hl{1327.62 万吨（同比 +16.68\%）}，
为历史第三高位\cite{yuntang2026sugar}。
需要说明的是，\hl{该榨季的产量口径在机构间存在显著分歧}：
中国糖业协会在 2025 年 11 月的预测为 \num{1170} 万吨，
而云糖网 2026 年 9 月的最终统计为 \num{1327.62} 万吨，两者相差 \num{157.6} 万吨
——\emph{榨季产量数据在榨季结束前具有很大的修正空间}，
这也是白糖价格在榨季中期出现预期修正式波动的原因之一。产量的连续回升来自糖料种植面积的恢复，
而种植面积的恢复又直接源于国内糖价在此前的上行（蔗农收益改善）——
这是政策定价型品种的典型传导：\emph{价格信号通过补贴与收购价传导到面积，再经过一个榨季传导到产量}。

需求侧的边际变化在国际市场：2025 年 1---10 月累计进口食糖 390.16 万吨（同比 +13.69\%），
糖浆与白砂糖预混粉进口约 100 万吨\cite{yuntang2026sugar}。
2025/26 年度全球食糖过剩量被国际糖业组织（ISO）从 220 万吨下调至 110 万吨，
并预计 2026/27 年度转为 \hl{20 万吨供应缺口}\cite{moa2026sugar}；
但机构间分歧极大：仅 2025/26 年度的过剩量估计，
就从 ISO 的 \num{110} 万吨（下调后）到 USDA 的 \num{1140} 万吨不等。
2026 年 8 月纽约原糖累计上涨超过 20\%、创 2010 年以来最大月度涨幅，
9 月 10 日 ICE 11 号原糖 10 月合约收于 18.76 美分/磅\cite{yuntang2026sugar}。

价格数据对这种“外强内弱”的结构给出了直接印证：2021 年 2 月至 2026 年 9 月，
郑商所白糖期货仅上涨 \hl{0.8\%}（图 \ref{fig:fut-norm-soft}），
而同期国际原糖经历了完整的“过剩$\to$短缺”切换。
国内价格被进口配额、国储与加工糖成本三重约束在一个窄区间内，
这正是白糖主窗口年化波动率仅 8.7\%、为全部农业期货中第二低的原因（图 \ref{fig:vol-bar}）。

\subsection{棉花：目标价格补贴与全球纺织需求}
\label{sec:soft-cotton}

棉花是中国农业中\hl{政策介入最深}的品种。2025 年全国棉花产量 664.1 万吨
（同比 +7.7\%），播种面积 2979.2 千公顷（+5.0\%），其中新疆产量 616.5 万吨
（占全国的 92.8\%）\cite{stats2025cotton}。供给的绝对集中意味着
\emph{新疆一地的天气就是全国供给的天气}。

政策方面，2026---2028 年新一轮新疆棉花目标价格政策延续每吨 18600 元的价格水平、
固定产量 510 万吨补贴，但有两处重要调整：\hl{质量补贴资金分配比例由 5\% 左右提高到 10\% 左右}，
\hl{长绒棉交售量补贴标准由细绒棉的 1.5 倍提高到 2 倍}\cite{xj2026cottonpolicy}。
政策取向从“保产量”转向“提质”——这与中央一号文件的种业振兴主线一致\cite{cccp2025doc1}。
进口方面，2025/26 年度累计进口棉花 158.5 万吨（同比 +50.7\%），
有关部门在 2025 年 8 月与 2026 年 3 月两次增发滑准税进口配额合计 50 万吨\cite{jinqiao2026cotton}。
\emph{配额的发放节奏本身构成一个政策周期}：配额宽松$\to$进口增加$\to$国内价格承压$\to$下一轮配额收紧。

从价格看，棉花期货在 2021---2026 年间上涨 3.5\%，但路径剧烈——先因 2021---2022 年
全球补库与新疆减产大涨至 21507 元/吨（主窗口最高），随后在纺织需求走弱中回落。
第 \ref{sec:corr-cycle} 节已经指出，棉花的周期分量与生猪的周期分量呈\hl{显著负相关（$-0.59$）}，
其根源正是两者驱动因素的完全错位：棉花由上一年度的全球纺织需求与新疆天气决定，
生猪由国内产能决定。

\subsection{苹果与红枣：真正的“产量周期”}
\label{sec:soft-fruit}

如果要在农业中寻找教科书式的“蛛网周期”，答案不是生猪，而是\hl{苹果}。

苹果的产量存在明确的\hl{“大小年”（隔年结果）}规律：连续两年产量相差 20\% 以上即认为出现
大小年现象，晚熟品种重于早中熟品种、老年树重于幼年树\cite{cma2015fruit}。
其机理是果树体内源激素与营养分配的自我调节——大年结果过多会消耗花芽分化所需的养分，
使次年成为小年。叠加花期霜冻与成熟期降水，形成“产量$\to$价格$\to$管理投入$\to$产量”的
短周期闭环。

数据清晰地展示了这一脉冲：全国苹果产量 2023/24 产季 3554.11 万吨\cite{fzzq2026apple}、
2024/25 产季升至 3734.97 万吨（+5.09\%）、
2025/26 产季回落至 \hl{3432.44 万吨（同比 -312 万吨，近七年最低）}\cite{fzzq2026apple}；
库存同步收缩，2025/26 产季全国入库峰值 735.77 万吨，为近五年最低、同比下降 104.57 万吨
\cite{fzzq2026apple}。产量与库存的双降推动苹果期货在 2026 年从低位反弹，
使其成为主窗口内唯一“逆势上涨”的种植类品种（+19.8\%，图 \ref{fig:fut-norm-soft}）。

红枣与苹果同属“产地气候 + 库存”驱动，但状态相反：
2024 产季新疆灰枣产量 75 万吨（较 2023 产季大增 67\%）\cite{mysteel2025jujube,qhrb2026jujube}，
2025 产季回落至 57.38 万吨（-23.49\%），但\emph{减产幅度有限且旧季库存高企}，
测算 2025/2026 产季结转库存将达 30 万$\sim$35 万吨、创近五年新高\cite{qhrb2026jujube}。
因此红枣呈现“减产不涨价”的形态：期货价格从主窗口初期的 15799 元/吨跌至 7579.7 元/吨
（\hl{-25.9\%}），主窗口年化波动率 25.9\%、仅次于生猪。
\hl{红枣是“库存周期压倒产量周期”的典型案例}——这一区分对第二部分的加工类企业尤其重要。

\subsection{花生：四年一轮回}
\label{sec:soft-peanut}

花生提供了一个“小品种、大波动”的完整样本。其需求结构为压榨 55\%$\sim$60\%、
食用 30\%$\sim$35\%、饲用 120 万$\sim$150 万吨，呈现“压榨定大盘、食用定品质”的双轨特征
\cite{nqw2026peanut}。由于压榨需求与油脂价格挂钩，花生在定价逻辑上
\hl{与豆油、棕榈油等油脂品种的联系强于与粮食的联系}。

供给端经历了一轮完整的四年周期：2022 年的短缺（价格高点）刺激种植面积扩张，
2024 年全国花生种植面积约 7256 万亩、产量 1898.3 万吨（连续第二年增产），
2025 年面积再增 4.01\% 至 457.49 万公顷（约 6862 万亩），为连增第三年\cite{nqw2026peanut}。
价格随之走弱，2025 年现货运行区间约 3.8---4.5 元/斤\cite{nqw2026peanut}。
期货端，花生主力合约从 2021 年的 10882.7 元/吨高位回落 24.2\%，
但主窗口年化波动率仅 12.1\%，是农产品期货中较低者——\emph{因为产量分散、
无单一产区天气风险，且价格由压榨利润锚定}。

\subsection{天然橡胶与纸浆：需求在外、供给在外}
\label{sec:soft-rubber}

这两个品种严格说来不属于“中国农业的周期”，而是\hl{中国农业资本参与全球资源品周期}的样本，
但它们在申万分类中对应种植业（天然橡胶）与轻工制造上游，且是农业上市公司
（如海南橡胶、部分林浆纸企业）的核心价格变量，因此纳入本章。

天然橡胶的全球供给已经连续多年偏紧：国际橡胶生产国协会（ANRPC）预计 2025 年全球天然橡胶
供应缺口将\hl{连续第五年存在}，全球产量仅增 0.3\% 至 1490 万吨、消费量增 1.8\% 至 1560 万吨；
印尼因农民转向油棕，2025 年产量预计同比下降 9.8\% 至 204 万吨，
越南减少 1.3\% 至 128 万吨\cite{anrpc2025rubber}。需求端，中国约 70\% 的天然橡胶用于轮胎生产
\cite{anrpc2025rubber}，因此橡胶周期本质上是\emph{汽车与物流周期}。
2021 年 2 月至 2026 年 9 月，天然橡胶上涨 \hl{24.8\%}、20 号胶上涨 43.3\%，
是主窗口内表现最好的农业相关品种之一。

纸浆的路径相反：2024 年主力合约波动区间 5500---6500 元/吨，2025 年四季度在进口浆报价
推动下反弹（12 月 31 日进口阔叶浆均价 4627.63 元/吨、较 9 月末 +9.08\%），
但 2026 年上半年转为偏弱下行，6 月末山东银星针叶浆现货 4720 元/吨（较年初 -880 元/吨），
港口库存维持 232 万吨以上的高位\cite{guoyuan2026pulp}。
结果是从 2021 年 2 月的高点回落 \hl{-29.8\%}，为本章跌幅最大的品种。
\hl{纸浆的周期是“浆厂投产 + 港口库存”的库存周期，与中国农产品供需无关}——
它在相关矩阵中相对孤立（图 \ref{fig:corr-heat}），正是这一独立性的统计表现。

\subsection{软商品的整体图景}
\label{sec:soft-summary}

\begin{figure}[htbp]
\centering
\begin{tikzpicture}
\begin{axis}[agri axis, yearticks,
  width=14.2cm, height=6.8cm,
  xlabel={年份（2021 年 2 月 $=100$）}, ylabel={归一化价格指数},
  xmin=2021.0, xmax=2026.9, ymin=55, ymax=165,
  xtick={2021,2022,2023,2024,2025,2026},
  legend style={at={(0.99,0.03)}, anchor=south east, legend columns=3,
                cells={anchor=west}, font=\tiny},
  legend entries={白糖, 棉花, 棉纱, 天然橡胶, 20 号胶, 苹果, 红枣, 纸浆, 花生},
]
\addplot[SoftRed, thick] table[x=x, y=SR]{data/clean/fut_norm_soft.dat};
\addplot[AquaTeal, thick] table[x=x, y=CF]{data/clean/fut_norm_soft.dat};
\addplot[AgriLight, thick] table[x=x, y=CY]{data/clean/fut_norm_soft.dat};
\addplot[InkGray, thick] table[x=x, y=RU]{data/clean/fut_norm_soft.dat};
\addplot[OilBlue, thick] table[x=x, y=NR]{data/clean/fut_norm_soft.dat};
\addplot[AgriGreen, thick] table[x=x, y=AP]{data/clean/fut_norm_soft.dat};
\addplot[AgriGold, thick] table[x=x, y=CJ]{data/clean/fut_norm_soft.dat};
\addplot[violet, thick] table[x=x, y=SP]{data/clean/fut_norm_soft.dat};
\addplot[Grain, thick] table[x=x, y=PK]{data/clean/fut_norm_soft.dat};
\node[font=\scriptsize, OilBlue, anchor=west] at (axis cs:2026.02,146) {20 号胶 143.3};
\node[font=\scriptsize, AgriGold, anchor=west] at (axis cs:2026.02,71) {红枣 74.1};
\node[font=\scriptsize, violet, anchor=west] at (axis cs:2026.02,66) {纸浆 70.2};
\end{axis}
\end{tikzpicture}
\caption{软商品与经济作物期货归一化走势（2021 年 2 月 $=100$）}
\label{fig:fut-norm-soft}
\srcfile{新浪财经期货主力连续合约；作者计算。}
\end{figure}

把本章与第 \ref{sec:grain} 章放在一起，可以得到软商品的三个总结论：

\begin{keybox}[软商品周期的三条规律]
\normalsize
\textbf{① 政策越深，波动越低。}白糖（8.7\%）、棉花（17.0\%）在配额、补贴与储备制度的
约束下波动受限；而苹果（20.4\%）、红枣（25.9\%）几乎无政策干预，波动完全由产量与库存释放。
\par\vspace{3pt}
\textbf{② 单一产区=单一风险。}新疆棉花（占全国 92.8\%）、新疆灰枣、
陕西与山东苹果（合计占全国一半以上），使这些品种的周期性表现为
“一地天气决定全国价格”的脉冲式波动，而非平滑周期。
\par\vspace{3pt}
\textbf{③ 库存决定“减产是否涨价”。}红枣在 2025 产季减产 23.5\% 的情况下价格仍下跌，
原因是结转库存创五年新高；苹果在减产 8.4\% 的同时库存降至五年最低，价格因此反弹。
\hl{产量周期必须与库存周期同时判断，这是种植类品种与养殖类品种最大的方法论差异}——
生猪无法库存，而果实可以。
\end{keybox}

\section{其他养殖与水产品：肉禽、蛋禽、牛羊、水产}
\label{sec:other-livestock}

生猪之外的养殖品种（肉牛、肉羊、白羽肉鸡、蛋鸡、水产）没有期货合约，
但它们构成了中国农业产值的另一大半，而且\hl{与猪周期共享同一条成本曲线、
同一种产能机制、同一批消费者}。本章的问题是：这些品种的周期是猪周期的“变奏”，
还是各自独立的周期？答案取决于一个变量——\emph{产能调整所需的时间}。

\subsection{产能时间轴：不同养殖品种的周期长度}
\label{sec:other-biology}

\begin{center}
\begin{tikzpicture}[font=\scriptsize, x=0.0082cm, y=1cm]
  % 时间轴：0 ~ 52 个月（200 刻度单位 = 6 个月）
  \draw[-{Stealth[length=4pt]}, InkGray, line width=0.8pt] (0,0) -- (1760,0);
  \foreach \m/\lb in {0/0, 200/6, 400/12, 600/18, 800/24, 1100/33, 1400/42}
    \draw[InkGray] (\m,0) -- (\m,-0.12) node[below, font=\tiny] {\lb};
  \node[font=\tiny, InkGray, anchor=west] at (1620,-0.30) {个月};
  % 各类品种的产能响应条
  \draw[AgriGold, line width=7pt, line cap=round] (0,0.70) -- (400,0.70);
  \node[font=\scriptsize, anchor=west, AgriGold!70!black] at (0,1.02) {蛋鸡、肉鸡：约 6---12 个月};
  \draw[HogPink, line width=7pt, line cap=round] (0,1.55) -- (550,1.55);
  \node[font=\scriptsize, anchor=west, HogPink!70!black] at (0,1.87) {生猪：约 10---12 个月};
  \draw[AquaTeal, line width=7pt, line cap=round] (0,2.40) -- (900,2.40);
  \node[font=\scriptsize, anchor=west, AquaTeal!70!black] at (0,2.72) {肉羊：约 18---24 个月};
  \draw[OilBlue, line width=7pt, line cap=round] (0,3.25) -- (1650,3.25);
  \node[font=\scriptsize, anchor=west, OilBlue!80!black] at (0,3.60) {肉牛：约 24---44 个月以上};
  % 参照线：生猪的产能周期终点
  \draw[InkGray, dashed, line width=0.5pt] (550,-0.1) -- (550,2.9);
\end{tikzpicture}
\end{center}

这张图解释了本章最核心的判断：\hl{产能时间越长的品种，价格周期越长、振幅越大、
与猪周期的错位越明显}。肉牛从犊牛到出栏约 24 个月，从补栏新生母牛到育成牛出栏
通常需 44 个月以上，且牛为单胎动物\cite{ndrc2026beef}——
这意味着 2023 年开始的母牛去化要到 2027 年之后才能反映为供给收缩；
而蛋鸡只需淘汰与补栏，去产能可在半年内完成。
因此当四个品种同时处于下行期时，它们的\emph{底部时间完全不同}。

\subsection{肉牛：政策触发的产能出清}
\label{sec:other-beef}

肉牛是本章中周期最完整、也最典型的外生冲击案例。价格自 2022 年末持续下降至 2025 年一季度，
2025 年二季度起止跌转涨，结束了持续两年的下行周期\cite{ndrc2026beef}。
产能方面，肉牛存栏于 2023 年三季度达到 \hl{10524 万头的峰值}，
随后持续下降，2025 年一季度降至 9462 万头（2020 年二季度以来最低），
2025 年末全国牛存栏 9608 万头（同比 -4.37\%），
2026 年二季度末进一步下滑至 9557 万头（同比 -4.3\%）\cite{ndrc2026beef,caas2026agri}。
同期出栏量与产量仍在增长（2025 年出栏 5133 万头、牛肉产量 801.00 万吨、创近五年新高），
说明\hl{当前的产量增长来自“产能去化后的集中出栏”}，是不可持续的。

价格已经开始反映供给收缩：2025 年肉牛交易均价 26.45 元/公斤（同比 +16.68\%），
国产鲜牛肉批发均价 59.53 元/公斤（+17.56\%）\cite{caas2026agri}；
2026 年 8 月末至 9 月初牛肉批发市场周均价 69.59 元/公斤（同比高 6.8\%）\cite{xinhua2026week}。
更关键的政策变量出现在 2025 年 12 月 31 日：商务部公告裁定进口牛肉数量增加
对中国国内产业造成严重损害，决定以“国别配额及配额外加征关税”形式实施保障措施，
期限 3 年（2026 年 1 月 1 日---2028 年 12 月 31 日），配额外加征关税税率 \hl{55\%}
\cite{mofcom2026beef}。这将进口牛肉（此前八件套均价约 48.64 元/公斤）
的价格优势大幅削弱，\emph{等于给国内肉牛价格加上了一个政策底}。

\subsection{肉羊：被猪价拖累的独立周期}
\label{sec:other-mutton}

肉羊提供了本章中最有趣的“交叉抑制”样本。产能上，经过两年多的调整，
2025 年第一季度末全国肉羊存栏 3.00 亿只，处于 2017 年以来本轮周期的低位；
2024 年新生羔羊数量同比下降 4.8\%，2025 年 1---5 月继续下降 3.36\%\cite{lj2026mutton}。
按产能逻辑，价格应当上涨，但 2025 年羊肉价格只走出“先抑后扬”的温和反弹：
1 月 70.1 元/kg 降至 7 月 68.9 元/kg，再回升至 12 月 71.6 元/kg，全年仅上涨 2.1\%
\cite{lj2026mutton}。

原因在于替代品的拖累：2025 年猪肉集贸市场均价同比下降 \hl{18.9\%}、
鸡肉价格下跌 5.8\%，作为羊肉核心替代品的牛肉价格虽上涨 8.5\% 但整体仍处近五年相对低位，
对羊肉价格形成明显的\hl{交叉抑制效应}\cite{lj2026mutton}。
到 2026 年，随着猪价企稳与牛羊供给收缩共振，羊肉批发价在 2026 年第 36 周
达到 66.13 元/公斤（同比高 8.9\%）\cite{xinhua2026week}。
\hl{肉羊的案例说明：产能周期决定方向，替代品价格决定幅度}。

\subsection{白羽肉鸡与蛋鸡：短周期、快反转}
\label{sec:other-poultry}

禽类的产能调整最快，因此周期最短、反转最急。白羽肉鸡在 2025 年经历了
“祖代更新创历史新高 + 出栏价格创近十年低位 + 产业链利润断崖”的组合：
祖代年度更新 157.41 万套（+4.89\%，超过 2013 年历史峰值），
其中海外引种占 39.92\%、国内自繁品种占 60.08\%；
在产父母代种鸡 5604.58 万套（+11.86\%），连续 3 年创历史新高；
全年出栏均价仅 7.28 元/kg（同比 -6.60\%，近 10 年低位）\cite{yangji2026chicken}。
出栏量却仍在增长：2025 年白羽肉鸡出栏 90.57 亿只（+5.97\%），
鸡肉产量 1930.19 万吨（+8.50\%），全国鸡肉产量 2803.44 万吨（+6.18\%）\cite{yangji2026chicken}。

反转发生在 2026 年：1---4 月毛鸡平均价格 3.63 元/斤（同比 +4.31\%），
毛鸡养殖利润 0.71 元/只（同比 +\hl{297.22\%}），
祖代引种因海外引种中断而全部转为国内自繁（更新量 29.38 万套、同比 -13.84\%）
\cite{yangji2026chicken}。\emph{引种中断是一个被低估的供给冲击}——
祖代供给的下滑将在 14 个月后传导到父母代、再传导到商品代。

蛋鸡的周期与生猪高度同步但幅度更大：2025 年全国在产蛋鸡年均存栏约 12.2 亿羽（+1.3\%），
9 月峰值突破 12.7 亿羽，全年平均淘汰日龄预计超过 520 天（去产能缓慢），
鸡蛋每千克平均亏损约 0.93 元\cite{aige2025layer}。
2026 年在产蛋鸡存栏逐月下降，8 月约 12.82 亿羽（同比 -6.08\%），
主产区蛋价于 2026 年 5 月底冲至 5 元/斤、创近十年同期新高\cite{zhuochuang2026layer}；
2026 年第 36 周鸡蛋批发价 10.83 元/公斤（同比 +\hl{41.9\%}）\cite{xinhua2026week}，
上半年禽蛋产量 1690 万吨（同比 -2.2\%）\cite{stats2026h1}。
在第 \ref{sec:corr-transmission} 节的传导检验中，\emph{鸡蛋是唯一稳健的、
与生猪周期同源的品种}，这与其“同样的产能机制 + 更短的调整时间”完全一致。

\subsection{水产：产能周期之外的另一个周期}
\label{sec:other-aqua}

水产品与前四类有本质区别：\hl{它的供给由养殖面积与投苗量决定，而不是由存栏决定}，
且需求端由居民蛋白消费结构升级驱动，与猪价的关系更弱
（第 \ref{sec:corr-sw} 节中水产养殖指数与猪价呈显著反周期）。

规模上，2025 年全国水产品总产量 7652.14 万吨（同比 +4.00\%），
其中养殖产量 6325.62 万吨（+4.38\%）、捕捞产量 1326.52 万吨（+2.23\%），
养殖与捕捞之比为 82.7:17.3\cite{moa2026fishery}；
水产品进口量 707.47 万吨（+2.05\%），贸易逆差扩大至 42.18 亿美元\cite{moa2026fishery}。

价格周期由 2024 年的\hl{投苗量大幅下降}启动：2025 年二季度起淡水鱼价格大幅上涨，
草鱼塘口均价 8.5 元/斤（同比 +\hl{43\%}）、价格高点达 14 元/公斤并创历史同期新高
\cite{nfnews2025grasscarp}；而海水产品同期走弱（南美白对虾 -8.98\%）\cite{caas2026agri}。
但到 2026 年第 36 周，淡水鱼价格全面回落：鲤鱼 14.12、草鱼 17.85、鲫鱼 20.22 元/公斤，
环比均小幅下跌\cite{xinhua2026week}。
\hl{这构成一个典型的“投苗$\to$上市”滞后周期（约 12---18 个月），
与生猪的“配种$\to$出栏”周期同构，但时间更长、弹性更大}。

\subsection{养殖业的周期矩阵}
\label{sec:other-matrix}

\begin{table}[htbp]
\centering
\small
\begin{tabular}{@{}lrrrrl@{}}\toprule
品种 & 产能周期 & 2025 年状态 & 2026 年方向 & 与猪周期 & 关键变量 \\
\midrule
生猪   & 10$\sim$12 月 & 深度亏损、政策去化 & 磨底回升 & 基准      & 能繁母猪、政策 \\
肉牛   & 24$\sim$44 月 & 存栏 -4.4\%、价格上涨 & 上行      & 错位      & 进口保障措施 \\
肉羊   & 18$\sim$24 月 & 存栏低位、价格微涨 & 温和上行  & 被抑制    & 替代品价格 \\
白羽肉鸡 & 6$\sim$12 月 & 价格近十年低位、利润崩塌 & 快速反转 & 高度同步  & 祖代引种 \\
蛋鸡   & 6$\sim$10 月 & 存栏高位、全面亏损 & 快速反转  & 高度同步  & 淘汰日龄 \\
水产   & 12$\sim$18 月 & 投苗下降、价格冲高 & 高位回落  & 反周期    & 投苗量、进口 \\
\bottomrule
\end{tabular}
\caption{养殖业各品种的周期特征对照（2025---2026）}
\label{tab:other-livestock-matrix}
\srcfile{数据来自本章引用的农业农村部、国家统计局、商务部公告与各产业技术体系报告。}
\end{table}

\begin{keybox}[本章三个结论]
\normalsize
\textbf{① 产能时间决定周期形态。}肉牛的 44 个月产能刚性使其供给缺口无法在短期内回补，
价格上行一旦形成就具有持续性；禽类的半年产能使其价格反复快速反转，
“高价$\to$亏损$\to$去化$\to$高价”的循环在两年内可以完成两轮。\par\vspace{3pt}
\textbf{② 替代效应是养殖业内部的传导机制。}2025 年猪价下跌 18.9\% 直接压制了羊肉价格，
使肉羊在自身产能已去化的条件下仅录得 2.1\% 的涨幅。\hl{养殖品种之间的相关性通过
“消费替代”而非“成本传导”实现}，这与第 \ref{sec:corr} 章商品价格相关性的结论形成互补。\par\vspace{3pt}
\textbf{③ 政策在养殖业中的作用从“托底”走向“出清”。}生猪产能调控方案、
进口牛肉保障措施、祖代引种中断，都是行政力量直接改变供给曲线的例子。
在 2025---2026 年的农业周期中，\hl{政策是产能去化的加速器，而不是价格的保证}。
\end{keybox}

\section{跨品种周期相关性与传导：以猪周期为基准}
\label{sec:corr}

本章回答本报告第一部分的核心问题：\hl{以猪周期为基准，农业各子行业之间到底有多强的
周期性关联？}为了回答这个问题，我们先固定方法口径，再分四个层次递进：
\emph{同期相关}（相关矩阵）$\to$ \emph{周期同步}（周期分量）$\to$
\emph{领先滞后}（互相关）$\to$ \emph{传导检验}（HAC 回归），
最后把结论落到\emph{股票市场}（申万子行业指数）与\emph{其他养殖品种}（牛羊、禽、水产）上。

\subsection{方法：为什么不能简单看相关系数}
\label{sec:corr-method}

农产品价格序列有很强的趋势性与结构性断点（例如 2018 年非洲猪瘟、2024 年玉米进口政策转向），
如果直接用价格\emph{水平}做相关，得到的往往是“共同的趋势”而不是“共同的周期”。
因此本章的所有计算遵循四条纪律：

\begin{enumerate}
  \item \textbf{同频率、同窗口}。所有相关系数在\emph{同一时间窗口}、\emph{同一频率}上计算，
  绝不在不同样本期上比较。主窗口取 \hl{2021 年 2 月---2026 年 8 月}
  （生猪期货 2021 年 1 月上市后的完整月份，共 67 个月度点），
  长窗口取 2015 年 2 月---2026 年 8 月（用于以生猪\emph{现货}价格指数为基准的检验）。
  \item \textbf{收益率口径为主}。基准口径是月度对数收益率
  $r_t=\ln(P_t/P_{t-1})$，并剔除换月造成的异常跳空（$|r|>35\%$ 视为拼接跳空）。
  \item \textbf{周期分量口径为辅}。用价格对 12 个月移动均线的偏离率刻画“周期位置”，
  该口径能捕捉同步性，但序列的持续性会使有效样本数远小于名义样本数，
  因此\hl{只作为定性参照，不用于显著性判断}。
  \item \textbf{显著性阈值}。主窗口 $n=67$ 时 5\% 显著性的临界值为 $|\rho|>0.24$；
  季度（3 个月）收益率口径 $n\approx 20$，临界值 $|\rho|>0.46$。
  传导回归使用 \hl{Newey--West (1987) HAC 稳健标准误，滞后阶 6}，
  以应对月度数据的自相关与异方差。
\end{enumerate}

\subsection{第一层：同期相关矩阵}
\label{sec:corr-matrix}

\begin{figure}[htbp]
\centering
\begin{tikzpicture}
\begin{axis}[
  width=11.2cm, height=10.6cm,
  xtick={0,1,...,21}, ytick={0,1,...,21},
  xticklabels={生猪,鸡蛋,玉米,淀粉,粳米,豆一,豆二,菜籽,花生,豆粕,菜粕,豆油,
               菜油,棕油,白糖,棉花,棉纱,橡胶,20号胶,苹果,红枣,纸浆},
  yticklabels={生猪,鸡蛋,玉米,淀粉,粳米,豆一,豆二,菜籽,花生,豆粕,菜粕,豆油,
               菜油,棕油,白糖,棉花,棉纱,橡胶,20号胶,苹果,红枣,纸浆},
  xticklabel style={rotate=90, anchor=east, font=\tiny},
  yticklabel style={font=\tiny},
  y dir=reverse,
  colormap={agriheat}{rgb255=(150,45,40) rgb255=(252,250,245) rgb255=(28,88,58)},
  colorbar horizontal,
  colorbar style={width=7.5cm, height=0.32cm, at={(0.5,-0.13)},
    anchor=north, xtick={-1,-0.5,0,0.5,1}, xticklabel style={font=\tiny},
    xlabel={\scriptsize 相关系数 $\rho$}, xlabel style={font=\scriptsize}},
  point meta min=-1, point meta max=1,
  axis on top, tick align=outside,
]
\addplot[matrix plot*, mesh/cols=22, point meta=explicit, shader=flat]
  table[x=i, y=j, z=rho]{data/clean/corr_m_short.dat};
\end{axis}
\end{tikzpicture}
\caption{农产品期货月度对数收益率相关矩阵（2021 年 2 月---2026 年 8 月，67 个月度观测）}
\label{fig:corr-heat}
\srcfile{新浪财经期货主力连续合约；作者计算。绿色为正相关、红色为负相关；
$|\rho|>0.24$ 为 5\% 显著性水平。}
\end{figure}

图 \ref{fig:corr-heat} 是最重要的单张图，它给出三条清晰的结论：

\begin{keybox}[相关性的三个层级]
\normalsize
\textbf{① 强相关只存在于“同一条产业链”内部。}深绿色的区块集中在两处：
\emph{油脂压榨链}（豆油---棕榈油---菜油---豆粕---豆二，$\rho$ 普遍在 $0.5\sim0.8$）
与\emph{棉花---棉纱}（同原料上下游，$\rho$ 接近 1）、\emph{玉米---玉米淀粉}
（$\rho\approx0.6$）、\emph{天然橡胶---20 号胶}。这些品种共享同一个供给来源或同一个加工链条，
因此相关性是“结构性”的、跨周期稳定的。\par\vspace{3pt}
\textbf{② 跨链条相关性普遍很弱。}养殖（生猪、鸡蛋）与粮食（玉米、粳米）之间、
养殖与软商品（白糖、棉花、橡胶）之间，$\rho$ 大多落在 $|0.1|$ 以内。\par\vspace{3pt}
\textbf{③ 与生猪相关性最高的不是玉米或豆粕，而是花生与粳米。}
这提示“同属农产品”并不构成相关性的理由，\emph{共同的驱动因子}才是。
\end{keybox}

% 由 scripts/make_tables.py 自动生成，请勿手工编辑
\begin{table}[htbp]
\centering
\small
\begin{adjustbox}{max width=\textwidth}
\begin{tabular}{@{}lrrrr@{}}
\toprule
类别对 & 品种对数 & 平均相关 & 最大相关 & 最小相关 \\
\midrule
压榨链 -- 压榨链 & 10 & 0.559 & 0.862 & 0.300 \\
软商品 -- 软商品 & 10 & 0.411 & 0.887 & 0.054 \\
油料 -- 压榨链 & 20 & 0.378 & 0.753 & 0.045 \\
粮食 -- 粮食 & 3 & 0.353 & 0.729 & 0.149 \\
粮食 -- 油料 & 12 & 0.299 & 0.573 & -0.134 \\
粮食 -- 压榨链 & 15 & 0.299 & 0.452 & 0.042 \\
油料 -- 油料 & 6 & 0.259 & 0.380 & 0.088 \\
养殖 -- 养殖 & 1 & 0.226 & 0.226 & 0.226 \\
压榨链 -- 其他 & 5 & 0.221 & 0.312 & 0.110 \\
压榨链 -- 软商品 & 25 & 0.195 & 0.565 & -0.220 \\
养殖 -- 油料 & 8 & 0.185 & 0.477 & -0.021 \\
油料 -- 其他 & 4 & 0.163 & 0.299 & 0.055 \\
养殖 -- 粮食 & 6 & 0.159 & 0.325 & -0.072 \\
油料 -- 林果 & 8 & 0.142 & 0.390 & -0.038 \\
软商品 -- 林果 & 10 & 0.134 & 0.311 & -0.027 \\
养殖 -- 压榨链 & 10 & 0.132 & 0.291 & -0.112 \\
压榨链 -- 林果 & 10 & 0.117 & 0.272 & -0.024 \\
油料 -- 软商品 & 20 & 0.114 & 0.372 & -0.139 \\
粮食 -- 软商品 & 15 & 0.081 & 0.386 & -0.176 \\
粮食 -- 林果 & 6 & 0.069 & 0.285 & -0.127 \\
林果 -- 林果 & 1 & 0.053 & 0.053 & 0.053 \\
养殖 -- 其他 & 2 & 0.052 & 0.096 & 0.009 \\
粮食 -- 其他 & 3 & 0.041 & 0.115 & -0.084 \\
养殖 -- 软商品 & 10 & 0.034 & 0.259 & -0.143 \\
养殖 -- 林果 & 4 & 0.032 & 0.127 & -0.018 \\
软商品 -- 其他 & 5 & 0.005 & 0.226 & -0.180 \\
林果 -- 其他 & 2 & -0.073 & -0.010 & -0.136 \\
\bottomrule
\end{tabular}
\end{adjustbox}
\caption{同一时间窗口（2021-02～2026-08）内大类之间的平均相关系数}
\label{tab:corr-group}
\end{table}


\begin{figure}[htbp]
\centering
\begin{tikzpicture}
\begin{axis}[agri axis, catbar,
  width=12.6cm, height=7.6cm,
  xlabel={与生猪期货月度收益的相关系数 $\rho$}, xmin=-0.3, xmax=0.55,
  yticklabels from table={data/clean/corr_with_hog.dat}{label},
  ytick=data,
  axis y line*=left,
]
\addplot[xbar, bar width=4.5pt, fill=AgriGreen!70, draw=AgriGreen]
  table[x=rho, y expr=\coordindex]{data/clean/corr_with_hog.dat};
\addplot[SoftRed, dashed, line width=0.8pt, domain=-0.3:0.55] {0};
\addplot[InkGray, dotted, line width=0.8pt, domain=-0.3:0.55] {0.24};
\addplot[InkGray, dotted, line width=0.8pt, domain=-0.3:0.55] {-0.24};
\end{axis}
\end{tikzpicture}
\caption{各期货品种与生猪期货的月度收益率相关系数（排序）}
\label{fig:corr-with-hog}
\srcfile{虚线为 0，点线为 5\% 显著性阈值 $\pm0.24$（$n=67$）。}
\end{figure}

% 由 scripts/make_tables.py 自动生成，请勿手工编辑
\begin{table}[htbp]
\centering
\small
\resizebox{\textwidth}{!}{%
\begin{tabular}{@{}llr@{}}
\toprule
品种 & 大类 & 与生猪期货的相关系数 \\
\midrule
花生 & 油料 & +0.477 \\
粳米 & 粮食 & +0.286 \\
豆一 & 油料 & +0.256 \\
鸡蛋 & 养殖 & +0.226 \\
豆油 & 压榨链 & +0.156 \\
棉纱 & 软商品 & -0.143 \\
豆粕 & 压榨链 & +0.128 \\
红枣 & 林果 & +0.127 \\
棉花 & 软商品 & -0.126 \\
菜油 & 压榨链 & +0.119 \\
菜粕 & 压榨链 & -0.112 \\
纸浆 & 其他 & +0.096 \\
玉米 & 粮食 & +0.091 \\
玉米淀粉 & 粮食 & +0.083 \\
豆二 & 油料 & +0.058 \\
天然橡胶 & 软商品 & +0.052 \\
白糖 & 软商品 & +0.026 \\
油菜籽 & 油料 & -0.021 \\
苹果 & 林果 & -0.018 \\
20号胶 & 软商品 & +0.014 \\
棕榈油 & 压榨链 & +0.003 \\
\bottomrule
\end{tabular}
}
\caption{各品种与生猪期货月度收益率的相关系数（2021-02～2026-08）}
\label{tab:corr-hog}
\srcfile{作者计算；样本 67 个月，5\% 显著性阈值 $|\rho|>0.24$。}
\end{table}


图 \ref{fig:corr-with-hog} 与表 \ref{tab:corr-hog} 给出了“以猪周期为基准”的直接答案：
\hl{没有任何一个农产品期货品种与生猪期货的月度收益率相关性达到 5\% 显著水平}。
最高的花生（$+0.48$）与粳米（$+0.29$）分别是小品种与政策定价品种，
其相关性更可能来自共同的宏观流动性因素而非产业传导；豆一（$+0.26$）
与鸡蛋（$+0.23$）处在临界值附近；而玉米（$+0.09$）、豆粕（$+0.13$）
——\emph{理论上与猪周期关联最强的两个饲料原料品种}——都不显著。

\subsection{第二层：周期同步性（周期分量口径）}
\label{sec:corr-cycle}

把价格换成“对 12 个月均线的偏离率”后，相关性结构发生了变化（图 \ref{fig:corr-cyc}）：
\hl{棉花（$-0.59$）与棉纱（$-0.66$）与猪周期呈现明显的反向同步}，
花生（$+0.52$）、纸浆（$+0.40$）同向，白糖（$-0.30$）反向，豆粕（$+0.28$）弱同向。

\begin{figure}[htbp]
\centering
\begin{tikzpicture}
\begin{axis}[agri axis, catbar,
  width=12.6cm, height=7.0cm,
  xlabel={与生猪价格周期分量的相关系数 $\rho$}, xmin=-0.75, xmax=0.6,
  yticklabels from table={data/clean/corr_cyc_hog.dat}{label},
  ytick=data,
]
\addplot[xbar, bar width=4.5pt, fill=OilBlue!70, draw=OilBlue]
  table[x=rho, y expr=\coordindex]{data/clean/corr_cyc_hog.dat};
\addplot[SoftRed, dashed, line width=0.8pt, domain=-0.75:0.6] {0};
\end{axis}
\end{tikzpicture}
\caption{各品种与生猪（现货价格指数）周期分量的相关系数}
\label{fig:corr-cyc}
\srcfile{周期分量 $=$ 月度价格 / 12 个月移动均线 $-1$；主窗口同前。
该口径序列持续性强，\emph{不用于显著性判断}。}
\end{figure}

棉花与生猪的\emph{反向}关系有清晰的经济解释：2021---2022 年是棉花大牛市
（新疆减产与全球补库）而猪价崩塌，2023---2024 年棉花回落而猪价反弹，
两者的周期错位来自完全不同的驱动因素（棉花看全球纺织需求与新疆天气，
生猪看国内产能）。
这恰恰强化了第 \ref{sec:corr-matrix} 节的结论：\hl{跨品种相关性的方向本身不稳定}，
不能作为交易或风险管理的稳定依据。

\subsection{第三层：领先滞后结构}
\label{sec:corr-leadlag}

\begin{figure}[htbp]
\centering
\begin{tikzpicture}
\begin{groupplot}[group style={group size=2 by 2, horizontal sep=1.1cm, vertical sep=1.25cm},
  agri axis, width=6.4cm, height=4.5cm,
  xlabel style={font=\scriptsize}, ylabel style={font=\scriptsize},
  tick label style={font=\tiny},
  xmin=-12, xmax=12, xtick={-12,-6,0,6,12},
  ymin=-0.9, ymax=0.9, ytick={-0.8,-0.4,0,0.4,0.8}, ylabel={$\rho$}]
\nextgroupplot[title={（a）豆粕}, xlabel={滞后月数 $k$}]
  \addplot[OilBlue, thick] table[x=lag, y=rho]{data/clean/leadlag2_M.dat};
  \addplot[SoftRed, dashed, line width=0.7pt, domain=-12:12] {0};
  \addplot[InkGray, dotted, line width=0.7pt, domain=-12:12] {0.24};
  \addplot[InkGray, dotted, line width=0.7pt, domain=-12:12] {-0.24};
\nextgroupplot[title={（b）鸡蛋}]
  \addplot[AgriGold, thick] table[x=lag, y=rho]{data/clean/leadlag2_JD.dat};
  \addplot[SoftRed, dashed, line width=0.7pt, domain=-12:12] {0};
  \addplot[InkGray, dotted, line width=0.7pt, domain=-12:12] {0.24};
  \addplot[InkGray, dotted, line width=0.7pt, domain=-12:12] {-0.24};
\nextgroupplot[title={（c）玉米淀粉}, xlabel={滞后月数 $k$}]
  \addplot[Grain, thick] table[x=lag, y=rho]{data/clean/leadlag2_CS.dat};
  \addplot[SoftRed, dashed, line width=0.7pt, domain=-12:12] {0};
  \addplot[InkGray, dotted, line width=0.7pt, domain=-12:12] {0.24};
  \addplot[InkGray, dotted, line width=0.7pt, domain=-12:12] {-0.24};
\nextgroupplot[title={（d）白糖}]
  \addplot[SoftRed, thick] table[x=lag, y=rho]{data/clean/leadlag2_SR.dat};
  \addplot[SoftRed, dashed, line width=0.7pt, domain=-12:12] {0};
  \addplot[InkGray, dotted, line width=0.7pt, domain=-12:12] {0.24};
  \addplot[InkGray, dotted, line width=0.7pt, domain=-12:12] {-0.24};
\end{groupplot}
\end{tikzpicture}
\caption{与生猪周期分量的领先滞后互相关（$k>0$ 表示该品种滞后猪周期 $k$ 个月）}
\label{fig:leadlag}
\srcfile{周期分量口径，$\pm12$ 个月，点线为 5\% 显著性阈值；仅展示相关系数绝对值较大的品种。}
\end{figure}

% 由 scripts/make_tables.py 自动生成，请勿手工编辑
\begin{table}[htbp]
\centering
\small
\resizebox{\textwidth}{!}{%
\begin{tabular}{@{}lllrrrl@{}}
\toprule
品种 & 代码 & 大类 & 同期相关 & 最大显著相关 & 滞后月 & 判定 \\
\midrule
棉花 & CF & 软商品 & -0.588 & -0.816 & 2.000 & 生猪领先 \\
棉纱 & CY & 软商品 & -0.656 & -0.804 & 2.000 & 生猪领先 \\
玉米淀粉 & CS & 粮食 & -0.167 & -0.715 & 3.000 & 生猪领先 \\
纸浆 & SP & 其他 & +0.396 & -0.679 & 9.000 & 生猪领先 \\
白糖 & SR & 软商品 & -0.305 & +0.641 & 9.000 & 生猪领先 \\
红枣 & CJ & 林果 & -0.266 & +0.613 & -9.000 & 生猪滞后 \\
苹果 & AP & 林果 & -0.417 & +0.603 & -7.000 & 生猪滞后 \\
豆二 & B & 油料 & +0.182 & -0.573 & 8.000 & 生猪领先 \\
菜粕 & RM & 压榨链 & -0.254 & -0.570 & 3.000 & 生猪领先 \\
花生 & PK & 油料 & +0.520 & +0.525 & -1.000 & 生猪滞后 \\
棕榈油 & P & 压榨链 & -0.273 & +0.521 & -5.000 & 生猪滞后 \\
豆一 & A & 油料 & -0.043 & -0.517 & 5.000 & 生猪领先 \\
油菜籽 & RS & 油料 & -0.166 & +0.470 & -6.000 & 生猪滞后 \\
豆油 & Y & 压榨链 & +0.063 & -0.458 & 8.000 & 生猪领先 \\
豆粕 & M & 压榨链 & +0.283 & -0.458 & 6.000 & 生猪领先 \\
玉米 & C & 粮食 & -0.030 & +0.362 & -4.000 & 生猪滞后 \\
菜油 & OI & 压榨链 & -0.040 & -0.341 & 5.000 & 生猪领先 \\
天然橡胶 & RU & 软商品 & -0.134 & +0.294 & -9.000 & 生猪滞后 \\
20号胶 & NR & 软商品 & -0.134 & +0.286 & -6.000 & 生猪滞后 \\
鸡蛋 & JD & 养殖 & +0.224 & \nadata & \nadata & 无显著领先滞后关系 \\
\bottomrule
\end{tabular}
}
\caption{以生猪周期分量为基准的领先滞后关系（月度，±12 个月）}
\label{tab:leadlag}
\srcfile{作者计算；比较对象为价格对 12 个月移动平均的偏离（周期分量），只报告通过 5\% 显著性的滞后阶；滞后月 $k>0$ 表示生猪领先 $k$ 个月。}
\end{table}


图 \ref{fig:leadlag} 的形状比数值更有信息量：四条互相关曲线都在
$k\in[-2,3]$ 的窄区间内出现主峰，随后迅速衰减到噪声水平。
这说明\hl{农产品价格之间的“领先滞后”关系基本上在季度以内}，
不存在可利用的、跨越半年以上的稳定传导链。对报告的结论而言，
这意味着：\emph{不能用猪价的拐点去推断 6 个月后玉米或豆粕的拐点}。

\subsection{第四层：传导检验（HAC 回归）}
\label{sec:corr-transmission}

把“相关性”升级为“条件传导”，检验模型为

\begin{equation}
  r^{i}_{t} = \alpha + \beta_{k}\, r^{\text{LH}}_{t-k} + \varepsilon_{t},
  \qquad k = 0,1,\dots,12 ,
  \label{eq:transmission}
\end{equation}

其中 $r^{\text{LH}}$ 为生猪期货月度收益率，$r^{i}$ 为目标品种收益率，
标准误为 Newey--West HAC（滞后 6 阶）。表 \ref{tab:transmission} 报告显著滞后阶数与最优滞后阶。

% 由 scripts/make_tables.py 自动生成，请勿手工编辑
\begin{table}[htbp]
\centering
\small
\resizebox{\textwidth}{!}{%
\begin{tabular}{@{}lllrrrrl@{}}
\toprule
品种 & 代码 & 大类 & 显著滞后阶数 & 最优滞后 $k$ & $\beta$ & $t$ 值 & 显著 $k$ 清单 \\
\midrule
白糖 & SR & 软商品 & 7 & 2.000 & -0.091 & -4.170 & 1,2,4,6,7,8,11 \\
棉花 & CF & 软商品 & 5 & 1.000 & -0.299 & -4.740 & 1,2,4,5,11 \\
玉米淀粉 & CS & 粮食 & 4 & 11.000 & +0.081 & 2.930 & 2,3,5,11 \\
棉纱 & CY & 软商品 & 4 & 1.000 & -0.202 & -4.390 & 1,2,4,5 \\
豆一 & A & 油料 & 3 & 2.000 & -0.129 & -3.380 & 0,2,5 \\
豆粕 & M & 压榨链 & 3 & 8.000 & -0.252 & -3.380 & 6,8,12 \\
花生 & PK & 油料 & 2 & 0.000 & +0.214 & 4.910 & 0,2 \\
20号胶 & NR & 软商品 & 2 & 2.000 & -0.106 & -2.200 & 2,7 \\
纸浆 & SP & 其他 & 2 & 8.000 & -0.199 & -2.400 & 8,9 \\
棕榈油 & P & 压榨链 & 2 & 8.000 & -0.215 & -2.530 & 2,8 \\
豆二 & B & 油料 & 2 & 8.000 & -0.288 & -3.500 & 6,8 \\
油菜籽 & RS & 油料 & 2 & 2.000 & -0.141 & -4.200 & 2,11 \\
鸡蛋 & JD & 养殖 & 1 & 0.000 & +0.191 & 2.390 & 0 \\
菜粕 & RM & 压榨链 & 1 & 12.000 & +0.193 & 2.280 & 12 \\
苹果 & AP & 林果 & 1 & 2.000 & -0.177 & -2.860 & 2 \\
豆油 & Y & 压榨链 & 1 & 8.000 & -0.213 & -3.940 & 8 \\
玉米 & C & 粮食 & 0 & \nadata & \nadata & \nadata & 无 \\
菜油 & OI & 压榨链 & 0 & \nadata & \nadata & \nadata & 无 \\
天然橡胶 & RU & 软商品 & 0 & \nadata & \nadata & \nadata & 无 \\
红枣 & CJ & 林果 & 0 & \nadata & \nadata & \nadata & 无 \\
\bottomrule
\end{tabular}
}
\caption{生猪价格向各品种的传导回归（月度收益，HAC 稳健标准误）}
\label{tab:transmission}
\srcfile{作者计算；$r^{i}_t=\alpha+\beta_k r^{\text{LH}}_{t-k}+\varepsilon_t$，Newey--West 滞后阶 6，样本 2021-02～2026-08。}
\end{table}


\begin{figure}[htbp]
\centering
\begin{tikzpicture}
\begin{groupplot}[group style={group size=2 by 1, horizontal sep=1.3cm},
  agri axis, width=6.6cm, height=4.6cm,
  xlabel style={font=\scriptsize}, ylabel style={font=\scriptsize},
  tick label style={font=\tiny},
  xmin=-0.5, xmax=12.5, xtick={0,3,6,9,12},
  ymin=-5, ymax=5, ytick={-4,-2,0,2,4}, ylabel={$t$ 值},
  legend style={font=\tiny, at={(0.02,0.03)}, anchor=south west}]
\nextgroupplot[title={（a）豆粕}]
  \addplot[OilBlue, thick, mark=*, mark size=1.2pt]
    table[x=k, y=t]{data/clean/trans_M.dat};
  \addplot[SoftRed, dashed, line width=0.7pt, domain=-0.5:12.5] {1.96};
  \addplot[SoftRed, dashed, line width=0.7pt, domain=-0.5:12.5] {-1.96};
\nextgroupplot[title={（b）鸡蛋}]
  \addplot[AgriGold, thick, mark=*, mark size=1.2pt]
    table[x=k, y=t]{data/clean/trans_JD.dat};
  \addplot[SoftRed, dashed, line width=0.7pt, domain=-0.5:12.5] {1.96};
  \addplot[SoftRed, dashed, line width=0.7pt, domain=-0.5:12.5] {-1.96};
\end{groupplot}
\end{tikzpicture}
\caption{生猪收益率对目标品种收益率的滞后传导：HAC $t$ 值随滞后阶的变化}
\label{fig:transmission}
\srcfile{模型 \eqref{eq:transmission}；虚线为 $\pm1.96$（5\% 显著性）。}
\end{figure}

结果具有高度的规律性：\hl{21 个目标品种中有 17 个至少存在一个显著滞后阶，
但显著系数都很小（$|\beta|<0.3$）且符号不一致}。
对饲料原料而言，豆粕在 $k=6,8,12$ 显著但系数为 \emph{负}（$k=8$ 时 $\beta=-0.25$，$t=-3.38$），
菜粕在 $k=12$ 显著为 \emph{正}（$\beta=+0.19$），二者方向相反、逻辑上无法同时成立，
说明这些显著性更可能来自样本内的偶然共动而非稳定的经济机制
（\emph{21 个品种 $\times$ 13 个滞后阶 $=273$ 次检验，纯随机下也会有约 14 次“显著”}）。

这一多重检验问题必须显式处理，否则会得出“存在传导”的错误结论。
本报告的处理是把显著性当作\emph{筛选}而非\emph{证明}：
只有同时满足“系数方向符合产业链逻辑”、“在季度收益率口径下仍然显著”、
“在分样本（2021---2023 / 2024---2026）中符号一致”三个条件的传导关系才被认定为稳健。
按此标准，稳健的候选只剩三个：\hl{鸡蛋（$+0.28$ / $+0.26$）}、
豆一（$+0.30$ / $+0.18$）与花生（$+0.56$ / $+0.17$），
其中只有鸡蛋有明确的产业逻辑——它与生猪同属“养殖产能周期”驱动的品种，
且是唯一在 HAC 同期回归中显著的品种（$\beta=0.19$，$t=2.39$）。
这正是第 \ref{sec:other-livestock} 节把养殖业作为独立周期组来处理的原因。

\subsection{第五层：阶段收益对比}
\label{sec:corr-phase}

相关性度量的是“同时变动”，而下文更关心“在猪价上涨/下跌阶段，其他品种表现如何”。
以生猪价格的 ZigZag 分段为状态变量，计算各品种在猪价上行段与下行段的
月均收益率，结果见图 \ref{fig:phase-scatter} 与表 \ref{tab:phase-futures}。

\begin{figure}[htbp]
\centering
\begin{tikzpicture}
\begin{axis}[agri axis,
  width=12.4cm, height=7.4cm,
  xlabel={猪价上行段：品种月均收益（\%）},
  ylabel={猪价下行段：品种月均收益（\%）},
  xmin=-4.2, xmax=2.6, ymin=-1.8, ymax=2.2,
  xtick={-4,-3,-2,-1,0,1,2}, ytick={-1.5,-1,-0.5,0,0.5,1,1.5,2},
  legend style={at={(0.99,0.03)}, anchor=south east, cells={anchor=west}},
  legend entries={各品种, 45$^\circ$ 线（两阶段同收益）},
]
\addplot[SoftRed, dashed, line width=0.9pt, domain=-4.2:2.6] {x};
\addplot[only marks, mark=*, mark size=2.2pt, AgriGreen,
  nodes near coords, point meta=explicit symbolic,
  every node near coord/.append style={font=\tiny, InkGray, anchor=west}]
  table[x=up, y=down, meta=label]{data/clean/phase_scatter.dat};
\end{axis}
\end{tikzpicture}
\caption{生猪价格上行段与下行段中，各期货品种的月均收益率}
\label{fig:phase-scatter}
\srcfile{横轴为猪价上行段（共 15 个月）内的品种月均收益，纵轴为下行段（共 47 个月）内的月均收益；
点位于 45$^\circ$ 线上方说明该品种在猪价下跌时表现更好。}
\end{figure}

\begin{figure}[htbp]
\centering
\begin{tikzpicture}
\begin{axis}[agri axis, catbar,
  width=12.4cm, height=8.2cm,
  xlabel={与猪价同向的月份占比（\%）}, xmin=35, xmax=72,
  yticklabels from table={data/clean/phase_same_fut.dat}{label},
  ytick=data,
]
\addplot[xbar, bar width=4.5pt, fill=AgriGold!75, draw=AgriGold]
  table[x=pct, y expr=\coordindex]{data/clean/phase_same_fut.dat};
\addplot[SoftRed, dashed, line width=0.8pt, domain=35:72] {50};
\end{axis}
\end{tikzpicture}
\caption{与猪价同向的月份占比（50\% 为无关系基准）}
\label{fig:phase-same}
\srcfile{同向 $=$ 品种月度收益率与生猪价格指数变动方向一致。}
\end{figure}

% 由 scripts/make_tables.py 自动生成，请勿手工编辑
\begin{table}[htbp]
\centering
\scriptsize
\begin{adjustbox}{max width=\textwidth}
\begin{tabular}{@{}lllrrrrrr@{}}
\toprule
品种 & 代码 & 大类 & 上涨段月均\% & 下跌段月均\% & 差值（pp） & 同向月占比\% & 上涨月数 & 下跌月数 \\
\midrule
生猪 & LH & 养殖 & +5.36 & -3.56 & +8.92 & \nodata & 15 & 47 \\
鸡蛋 & JD & 养殖 & +0.53 & -0.69 & +1.22 & 64.50 & 15 & 47 \\
玉米 & C & 粮食 & -0.10 & -0.32 & +0.22 & 56.50 & 15 & 47 \\
玉米淀粉 & CS & 粮食 & -0.34 & -0.23 & -0.11 & 53.20 & 15 & 47 \\
粳米 & RR & 粮食 & -0.21 & +0.02 & -0.23 & 58.10 & 15 & 47 \\
强麦 & WH & 粮食 & +0.89 & +0.78 & +0.11 & 52.40 & 7 & 14 \\
豆一 & A & 油料 & -0.69 & -0.19 & -0.50 & 59.70 & 15 & 47 \\
豆二 & B & 油料 & -0.70 & -0.06 & -0.64 & 53.20 & 15 & 47 \\
油菜籽 & RS & 油料 & -0.69 & +0.21 & -0.90 & 50.00 & 15 & 47 \\
花生 & PK & 油料 & +1.81 & -1.09 & +2.91 & 67.20 & 14 & 47 \\
豆粕 & M & 压榨链 & -0.50 & -0.23 & -0.27 & 53.20 & 15 & 47 \\
菜粕 & RM & 压榨链 & -1.79 & +0.13 & -1.92 & 54.80 & 15 & 47 \\
豆油 & Y & 压榨链 & -0.45 & +0.34 & -0.79 & 59.70 & 15 & 47 \\
菜油 & OI & 压榨链 & -0.67 & +0.18 & -0.85 & 64.50 & 15 & 47 \\
棕榈油 & P & 压榨链 & -1.26 & +1.12 & -2.38 & 53.20 & 15 & 47 \\
白糖 & SR & 软商品 & -0.28 & +0.13 & -0.41 & 48.40 & 15 & 47 \\
棉花 & CF & 软商品 & -2.84 & +0.93 & -3.76 & 43.50 & 15 & 47 \\
棉纱 & CY & 软商品 & -1.90 & +0.55 & -2.45 & 41.90 & 15 & 47 \\
天然橡胶 & RU & 软商品 & +0.10 & +0.27 & -0.17 & 43.50 & 15 & 47 \\
20号胶 & NR & 软商品 & +0.06 & +0.47 & -0.41 & 48.40 & 15 & 47 \\
苹果 & AP & 林果 & -1.85 & +1.66 & -3.51 & 45.20 & 15 & 47 \\
红枣 & CJ & 林果 & -3.34 & +0.77 & -4.12 & 58.10 & 15 & 47 \\
纸浆 & SP & 其他 & +1.06 & -0.72 & +1.79 & 50.00 & 15 & 47 \\
\bottomrule
\end{tabular}
\end{adjustbox}
\caption{生猪期货上涨段与下跌段中各品种的月均收益（2021—2026）}
\label{tab:phase-futures}
\end{table}


散点图给出了一个反直觉但极其重要的事实：\hl{大多数品种位于 45$^\circ$ 线上方}，
即\emph{它们在猪价下跌阶段的表现好于猪价上涨阶段}。
这再次印证了第 \ref{sec:grain-cycle} 节的“量与价分离”：猪价上行往往对应
养殖产能扩张（利多饲料原料的\emph{数量}），但当猪价上行由供给冲击驱动时
（如非洲猪瘟），饲料需求反而是下降的；而猪价下跌阶段常伴随成本端的原料上涨
（2021---2022 年的玉米、豆粕牛市）。同向占比最高的花生（67.2\%）与鸡蛋（64.5\%）
也印证了显著性检验的结果。

\subsection{股票市场：申万农业子行业指数}
\label{sec:corr-sw}

对第一部分而言，股票市场提供了一个额外维度的验证。用申万一级（801010 农林牧渔）
与全部二级、三级农业行业指数的月度收益率与生猪期货收益率做相关，结果见图 \ref{fig:sw-corr}。

\begin{figure}[htbp]
\centering
\begin{tikzpicture}
\begin{axis}[agri axis, catbar,
  width=12.0cm, height=6.6cm,
  xlabel={与生猪期货月度收益的相关系数 $\rho$}, xmin=-0.25, xmax=0.06,
  yticklabels from table={data/clean/sw_corr_hog.dat}{label},
  ytick=data,
]
\addplot[xbar, bar width=4.5pt, fill=HogPink!75, draw=HogPink]
  table[x=rho, y expr=\coordindex]{data/clean/sw_corr_hog.dat};
\addplot[SoftRed, dashed, line width=0.8pt, domain=-0.25:0.06] {0};
\addplot[InkGray, dotted, line width=0.8pt, domain=-0.25:0.06] {-0.24};
\end{axis}
\end{tikzpicture}
\caption{申万农业子行业指数与生猪期货的月度收益率相关系数}
\label{fig:sw-corr}
\srcfile{申万行业指数（2021 版分类）月度收益率，主窗口 67 个月度点；未达 5\% 显著性。}
\end{figure}

% 由 scripts/make_tables.py 自动生成，请勿手工编辑
\begin{table}[htbp]
\centering
\scriptsize
\begin{adjustbox}{max width=\textwidth}
\begin{tabular}{@{}lllrrr@{}}
\toprule
行业代码 & 行业名称 & 层级 & 上涨段月均\% & 下跌段月均\% & 差值（pp） \\
\midrule
801\,010 & 农林牧渔 & 一级 & -1.00 & -0.41 & -0.59 \\
801\,012 & 农产品加工 & 二级 & -3.08 & +0.26 & -3.34 \\
801\,014 & 饲料 & 二级 & -0.85 & -0.95 & +0.10 \\
801\,015 & 渔业 & 二级 & -3.76 & +1.27 & -5.03 \\
801\,016 & 种植业 & 二级 & -2.22 & +0.52 & -2.74 \\
801\,017 & 养殖业 & 二级 & +0.50 & -0.96 & +1.47 \\
801\,018 & 动物保健 & 二级 & -4.62 & +0.60 & -5.22 \\
850\,111 & 种子 & 三级 & -2.88 & +0.27 & -3.15 \\
850\,113 & 其他种植业 & 三级 & -1.62 & +0.95 & -2.57 \\
850\,122 & 水产养殖 & 三级 & -4.23 & +1.32 & -5.55 \\
850\,142 & 畜禽饲料 & 三级 & -2.90 & -0.79 & -2.11 \\
850\,151 & 果蔬加工 & 三级 & -1.35 & +1.84 & -3.19 \\
850\,152 & 粮油加工 & 三级 & -4.57 & -0.89 & -3.69 \\
850\,154 & 其他农产品加工 & 三级 & -2.33 & +0.73 & -3.07 \\
850\,172 & 生猪养殖 & 三级 & -0.28 & -0.62 & +0.35 \\
850\,173 & 肉鸡养殖 & 三级 & -1.59 & -0.41 & -1.18 \\
850\,181 & 动物保健Ⅲ & 三级 & -4.62 & +0.60 & -5.22 \\
\bottomrule
\end{tabular}
\end{adjustbox}
\caption{申万农业子行业指数在生猪上涨段与下跌段的月均收益（2021—2026）}
\label{tab:phase-sw}
\end{table}


结论非常明确：\hl{申万农业各子行业指数与猪价的月度相关性全部在 $[-0.20,+0.03]$ 之间，
无一显著}，包括理论上关联最紧密的生猪养殖（$-0.15$）与畜禽饲料（$-0.13$）。
原因是股票价格交易的是\emph{预期}\emph{与估值}，而猪价是\emph{当期实现}的价格：
在猪价上行期，市场往往已经在交易产能扩张带来的未来供给过剩（股价见顶）；
在猪价下行期，市场又开始交易产能出清带来的反转预期。
表 \ref{tab:phase-sw} 的阶段收益对比印证了这一点，并给出一个清晰的分组：

\begin{itemize}
  \item \textbf{顺周期组（仅养殖）}：\emph{养殖业指数在猪价上行段月均 $+0.50\%$、
  下行段 $-0.96\%$（差值 $+1.47$ 个百分点）}，生猪养殖指数差值 $+0.35$ 个百分点；
  这是全部申万农业行业中唯一呈现“猪价涨则强、猪价跌则弱”的行业。
  \item \textbf{反周期组（幅度最大）}：水产养殖（$-5.55$ 个百分点）、
  动物保健（$-5.22$）、渔业（$-5.03$）、粮油加工（$-3.69$）。
  这些行业在猪价下跌阶段显著跑赢——水产养殖与渔业的需求由\emph{鱼类}
  而非生猪决定，动物保健则受益于猪价低迷期的防疫投入与出栏量增长。
  \item \textbf{中性组}：饲料（$+0.10$）、农林牧渔整体（$-0.59$）、
  其他农产品加工（$-3.07$）介于两者之间。
\end{itemize}

这一分组对第二部分有一个直接含义：\hl{上市公司画像不能按“农业”整体来做，
必须按产业链位置分别处理}——同为养殖业，生猪养殖与水产养殖的收益驱动完全相反。

\subsection{本章结论}
\label{sec:corr-conclusion}

\begin{keybox}[六个可检验的结论]
\normalsize
\textbf{① 相关性有层级，没有“农业整体周期”。}强相关只存在于同一产业链内部
（压榨链、棉纺链、玉米---淀粉、橡胶双品种）；跨链条相关性几乎全部不显著。
\par\vspace{3pt}
\textbf{② 与猪周期相关性最高的品种是花生与粳米，而非玉米、豆粕。}
理论上的“饲料传导链”在价格数据上不可见。
\par\vspace{3pt}
\textbf{③ 周期分量口径下才出现的强相关（棉花 $-0.59$、棉纱 $-0.66$）方向不稳定，}
只能作为定性参考。
\par\vspace{3pt}
\textbf{④ 领先滞后关系集中在 $\pm3$ 个月以内}，不存在可用的半年以上传导链。
\par\vspace{3pt}
\textbf{⑤ 传导回归的显著性大量来自多重检验；筛选后唯一稳健的候选是鸡蛋}
（养殖产能周期同源）。
\par\vspace{3pt}
\textbf{⑥ 股票市场对猪周期几乎没有线性定价}（全部 $|\rho|<0.20$），
只有养殖业指数显示出顺周期特征。
\end{keybox}

这些结论共同指向一个判断：\hl{中国农业的周期性是“分品种的、局部的”，
而不是“整体的、同一节拍”的}。种植业、养殖业、软商品各自的周期由不同的
外生变量驱动——养殖业看产能与疫病，粮食看政策与进口，软商品看全球供需与天气，
林业果业看产地气候与库存。因此，任何“农业整体周期”的叙述都必须先说明
它指的是哪一条链、哪一个环节。


% ---------------------------------------------------------------- 参考文献
\newpage
\printbibliography[heading=bibintoc, title={参考文献}]

\end{document}
